% STATUS: reviewed
\chapter{Modelių sluoksnis}\label{ch:models}

Kiekvienas kitas šios dalies skyrius aprašo vieną kelią, kuriuo modelio kas nors klausiama. Šis
skyrius aprašo tai, kas visiems tiems keliams bendra: kaip apskritai pasiekiamas modelis, kaip
atrodo vienas iškvietimas tinkle, kaip užtikrinama atsakymo forma, kas apie iškvietimą
užrašoma, kiek jis kainuoja ir kas už jį moka. Skyrius eina iš išorės į vidų.
\Cref{sec:models:door} nusako vienų durų į modelius principą; \cref{sec:models:providers} --
tiekėjus ir modelių identifikatorius; \cref{sec:models:generate,sec:models:embed} -- du
asistento paslaugos maršrutus, vieną generuojamam tekstui, kitą vektoriams;
\cref{sec:models:client} -- klientą, kuriuo API juos kviečia, ir jo klaidų apdorojimą. Antroji
pusė skirta pinigams ir limitams: kas užrašoma apie iškvietimą (\cref{sec:models:usage}),
biudžetai (\cref{sec:models:budgets}), tiekėjo instrukcijų podėlis ir kaip jis paveikė
sprendimus (\cref{sec:models:cache}), dažnio limitai (\cref{sec:models:ratelimits}) ir
modelis, kuriuo pagal numatymą veikia kiekvienas etapas, kartu su išmatuota pasirinkimo
priežastimi (\cref{sec:models:defaults}).


\section{Vienos durys į modelius}\label{sec:models:door}

Sistema veikia kaip kelios paslaugos (\cref{ch:arch}). Čia svarbios dvi iš jų.
\term{API} -- tai Django paslauga, kuri saugo visus duomenis ir per savo \term{užduočių
vykdytoją} atlieka fonines \term{užduotis}: perskaito teisės aktą, indeksuoja dokumentą,
parengia ataskaitą. \term{Asistento paslauga} (\code{legalops-assistant}: Node~24,
Express~5 ir AI~SDK 7~versija, kaip jas deklaruoja \file{package.json} kataloge
\file{assistant/app/}) yra vienintelis procesas, kuris apskritai jungiasi prie modelio. API
modelio nekviečia: kai užduoties etapui jo prireikia, API siunčia HTTP užklausą asistento
paslaugai, o ši atlieka iškvietimą ir grąžina atsakymą. Pokalbių asistentas veikia asistento
paslaugoje ir modelius pasiekia tuo pačiu kodu (\cref{ch:chat}).

Toks padalijimas sąmoningas, ir \file{docs/markdown/en/architecture.md} jį įtvirtina kaip dvi
projektavimo taisykles: \enquote{Django niekada nekviečia modelio tiekėjo} ir
\enquote{asistentas neturi jokių duomenų}. Iš to išplaukia trys dalykai.

\begin{itemize}
  \item \textbf{Tiekėjai ir raktai konfigūruojami vienoje vietoje.} Kokie tiekėjai yra, kokie
    jų adresai ir raktai ir kaip su kiekvienu kalbamasi, parašyta viename vienos paslaugos
    faile (\file{llm/provider.js}). Pridėjus tiekėją ar pakeitus raktą, daugiau nieko
    nekeičiama.
  \item \textbf{API niekada neturi rakto.} API parametruose nurodyta asistento paslauga
    (\code{ASSISTANT_URL}), jų tarpusavio durų bendroji paslaptis ir vektorizavimo modelio
    identifikatorius -- jokio tiekėjo adreso ir jokio tiekėjo rakto. Procesas, kuris atpažįsta
    įkeltų failų struktūrą ir aptarnauja viešuosius maršrutus, negali nutekinti to, ko neturi.
  \item \textbf{Kiekvienas iškvietimas registruojamas vienodai.} Kad ir kas klausė -- pokalbio
    ėjimas, užduoties etapas, vektorizavimo paketas, bandomasis paleidimas instrukcijų
    redaktoriuje, -- rezultatas yra viena eilutė lentelėje \code{llm_calls} su tais pačiais
    stulpeliais (\cref{sec:models:usage}). Kainos, biudžetai ir sąnaudų puslapis skaito tą
    vieną lentelę.
\end{itemize}

Asistento paslauga savo ruožtu nieko nesaugo: instrukcijos, pokalbiai ir sąnaudų eilutės
laikomos API duomenų bazėje, o asistento paslauga jas skaito ir rašo per API.

\paragraph{Vidinės durys.} Maršrutai, kuriuos kviečia API, yra po \code{/internal/llm/}
(\file{routes/internal.js}; \cref{tab:models:routes}). Jie egzistuoja tik izoliuotame tinkle
tarp konteinerių: viešasis atvirkštinis tarpinis serveris paslaugoms perduoda tik kelius po
\code{/api/}, todėl kelias, prasidedantis \code{/internal/}, iš išorės jų niekada nepasiekia.
Be to, kiekvienoje užklausoje turi būti antraštė \code{X-Internal-Secret}. Sargas
(\file{middleware/internal.js}, \code{requireInternal}) palygina ją su sukonfigūruota
paslaptimi per pastovų laiką ir tai daro dar prieš perskaitydamas užklausos turinį; jei
sukonfigūruota paslaptis tuščia, atmetamas kiekvienas iškvietimas su
403~\code{INTERNAL_ONLY}, užuot praleidus viską. API savo vidinius maršrutus saugo taip pat
(\file{assistant/api/internal_views.py}, \code{require_internal}). Užklausos turinys pro šias
duris gali būti iki 16~MB; viešasis pokalbio maršrutas, priimantis vieną žinutę, leidžia
512~kB.

\begin{table}[H]
\caption{Asistento paslaugos maršrutai, kuriuos kviečia API}\label{tab:models:routes}
\begin{tabularx}{\textwidth}{@{}L{58mm}Y@{}}
\toprule
\hd{Maršrutas} & \hd{Ką daro} \\
\midrule
\route{POST}{/internal/llm/generate} & Vienas generavimas: paprastas tekstas arba objektas, atitinkantis JSON schemą (\cref{sec:models:generate}). \\
\route{POST}{/internal/llm/embed} & Tekstų paketo vektoriai (\cref{sec:models:embed}). \\
\route{GET}{/internal/llm/models} & Sukonfigūruoti tiekėjai ir DI tinklų sąsajos modeliai su kainomis; kainoraščio (\cref{sec:models:usage}) ir instrukcijų redaktoriaus modelio pasirinkimo šaltinis. \\
\bottomrule
\end{tabularx}
\end{table}

\Cref{fig:models:path} rodo vieną iškvietimą nuo pradžios iki galo. Likusi skyriaus dalis
eina juo.

\begin{figure}[H]
\centering
\begin{tikzpicture}[x=1mm, y=1mm]
  % the API
  \node[svc, text width=27mm] (step) at (14,0) {užduoties etapas\\{\scriptsize arba maršrutas, kurio laukia žmogus}};
  \node[svc, text width=27mm] (rp) at (14,-18) {\code{run_prompt}\\{\scriptsize randa instrukcijos versiją, užpildo jos kintamuosius}};
  \node[svc, text width=27mm] (gen) at (14,-37) {\code{gateway.generate}\\{\scriptsize kartoja; išlaukia dažnio limito}};
  \node[chk, text width=27mm] (rec) at (14,-56) {\code{record_call}\\{\scriptsize \code{cost_of}: kiekviena įvesties dalis pagal savo kainą}};
  \node[store] (calls) at (14,-74) {\code{llm_calls}};
  \node[store] (pv) at (49,-18) {instrukcijų\\versijos};
  \node[store] (prices) at (49,-56) {\code{model_prices}};
  % the assistant service
  \node[chk, text width=28mm] (gate) at (89,-37) {\code{requireInternal}\\{\scriptsize bendroji paslaptis}};
  \node[svc, text width=28mm] (route) at (89,-56) {generavimo maršrutas\\{\scriptsize prašo srauto, perskaito jį iki galo}};
  \node[chk, text width=28mm] (check) at (89,-74) {objektas tikrinamas pagal schemą};
  % outside
  \node[ext, text width=17mm] (gw) at (130,-56) {DI tinklų sąsaja};
  \node[ai, text width=17mm] (model) at (130,-74) {modelis};

  \draw[flow] (step) -- (rp);
  \draw[flow] (rp) -- (gen);
  \draw[flow] (pv) -- node[lbl, above] {aktyvi} (rp);
  \draw[call] (gen) -- node[lbl, above] {POST \code{/internal/llm/generate}} (gate);
  \draw[flow] (gate) -- (route);
  \draw[call] (route) -- node[lbl, above] {Responses API} (gw);
  \draw[call] (gw) -- (model);
  \draw[flow] (route) -- (check);
  \draw[flow] (gen) -- node[lbl, right, pos=.45] {atsakymas ir penki žetonų skaičiai} (rec);
  \draw[flow] (prices) -- node[lbl, above] {kainos} (rec);
  \draw[flow] (rec) -- node[lbl, right] {viena eilutė} (calls);

  \begin{scope}[on background layer]
    \node[lane, fit=(step)(calls)(pv)(prices)] (laneapi) {};
    \node[lane, fit=(gate)(route)(check)] (laneas) {};
  \end{scope}
  \node[lanelabel, anchor=south west] at (laneapi.north west) {API (Django): užduočių vykdytojas arba žiniatinklio procesas};
  \node[lanelabel, anchor=south west] at (laneas.north west) {Asistento paslauga};
\end{tikzpicture}
\actorlegend
\caption{Vieno modelio iškvietimo kelias. Užduoties etapas kviečia \code{run_prompt}, kuri
instrukcijas ir modelį paima iš aktyvios instrukcijos versijos; API klientas siunčia užklausą
pro vidines duris; asistento paslauga prašo DI tinklų sąsajos srauto, perskaito jį iki galo,
patikrina objektą ir atsako visu atsakymu. Sąnaudos užrašomos vienoje vietoje: grįžus
atsakymui, \code{record_call} apskaičiuoja iškvietimo kainą ir įterpia vieną eilutę į
\code{llm_calls}. Punktyrinė rodyklė -- užklausa kartu su jos atsakymu.}\label{fig:models:path}
\end{figure}


\section{Tiekėjai ir modelių identifikatoriai}\label{sec:models:providers}

\term{Tiekėjas} -- tas, kas teikia modelius per API. Visoje sistemoje modelis vadinamas viena
\texttt{provider:model} formos eilute, pvz., \texttt{gateway:azure/gpt-6.1-sol}: priešdėlis,
parenkantis tiekėją, dvitaškis ir identifikatorius, kurį modeliui suteikė pats tiekėjas. Šią
eilutę instrukcijos versija saugo savo lauke \code{model} (\cref{ch:prompts}), ją užrašo
sąnaudų eilutė ir pagal ją sudarytas kainoraštis.

\subsection{Registras}

\file{llm/provider.js} laiko registrą: po vieną įrašą kiekvienam priešdėliui, nurodantį, ar
tiekėjas turi prisijungimo duomenis, kaip kuriamas jo klientas (vieną kartą, pirmą kartą
prireikus) ir kaip jis sukuria kalbos modelį bei vektorizavimo modelį
(\cref{tab:models:providers}). \code{parseModelId} eilutę skaido ties pirmuoju dvitaškiu, bet
tik tada, kai prieš dvitaškį stovi žinomas priešdėlis; bet kuri kita eilutė skaitoma kaip DI
tinklų sąsajos identifikatorius, todėl \model{azure/gpt-6.1-sol} ir
\texttt{gateway:azure/gpt-6.1-sol} yra tas pats modelis. Iškvietimo atsakyme modelis visada
nurodomas pilna forma (\code{qualifiedModel}), nes pagal ją API apskaičiuoja iškvietimo kainą.

\begin{table}[H]
\caption{Asistento paslaugos tiekėjų registras}\label{tab:models:providers}
\begin{tabularx}{\textwidth}{@{}L{21mm}L{37mm}YL{28mm}@{}}
\toprule
\hd{Priešdėlis} & \hd{Sukonfigūruotas, kai} & \hd{Kalbos modeliai} & \hd{Vektoriai} \\
\midrule
\texttt{gateway:} & nustatytas \code{AI_GATEWAY_KEY} (adresas: \code{AI_GATEWAY_URL}) & DI tinklų sąsajos Responses API, kiekvienas iškvietimas be būsenos & DI tinklų sąsajos \code{/embeddings} \\
\texttt{openai:} & nustatytas \code{OPENAI_API_KEY} & OpenAI, per atskirą SDK klientą & taip \\
\texttt{anthropic:} & nustatytas \code{ANTHROPIC_API_KEY} & Anthropic, per atskirą SDK klientą & nėra: jų paprašius -- klaida \\
\texttt{ollama:} & nustatytas \code{OLLAMA_BASE_URL} & Ollama serveris, kaip su OpenAI suderinami pokalbių modeliai & taip \\
\bottomrule
\end{tabularx}
\end{table}

Tiekėjas be prisijungimo duomenų sukelia aiškią klaidą: \code{providerFor} meta išimtį, ir maršrutas
atsako 400~\code{UNKNOWN_MODEL} dar prieš atliekant bet kokį iškvietimą. Iškvietimas niekada
neperduodamas kitam tiekėjui vietoje šio, o eilutė, nenurodanti jokio modelio, atmetama taip
pat, užuot ją persiuntus toliau, kad iškvietimas nepavyktų srauto viduryje.

Iš keturių šis diegimas naudoja vieną. \file{docs/markdown/en/architecture.md} užrašo sprendimą --
fakulteto DI tinklų sąsaja (angl. \en{gateway}) yra vienintelis modelių tiekėjas, -- o tinklas
jį užtikrina: vienintelis modelių serveris išeinančių jungčių leidžiamų adresų sąraše yra DI
tinklų sąsajos (\cref{ch:security}). Kiti trys priešdėliai neprivalomi: kiekvienas įjungiamas
nustačius jo raktą ar adresą, o šio diegimo compose failas asistento paslaugai neperduoda nė
vieno iš jų.

\subsection{DI tinklų sąsaja ir jos tinklo formatas}

\term{DI tinklų sąsaja} -- tai fakulteto Bifrost sąsaja (numatytasis adresas
\texttt{https://ai.knf.vu.lt/v1}). Ji naudoja OpenAI tinklo formatus ir kiekvieną iškvietimą
perduoda už jos įdiegtam modeliui; jos modelių identifikatoriai prasideda \code{azure/}.
Raktas siunčiamas iš karto dviejose antraštėse: kaip \code{x-bf-vk} ir kaip prieigos raktas
antraštėje \code{Authorization}.

Registras tuo pačiu raktu sukuria du klientus, nes dvi iškvietimų rūšys naudoja du skirtingus
formatus:

\begin{itemize}
  \item \textbf{Kalbos modeliai kviečiami per Responses API} (\code{/v1/responses}), o ne
    senesniu chat-completions formatu. Priežastis nurodyta \file{llm/provider.js}
    antraštiniame komentare: naujesni DI tinklų sąsajos modeliai (GPT-6 Astra, GPT-6.1 Sol)
    funkcijų įrankius priima tik ten, o pokalbių asistentas savo įrankius siunčia visada.
  \item \textbf{Vektoriai gaunami per} \code{/v1/embeddings}, naudojant su OpenAI suderinamą
    SDK klientą.
\end{itemize}

Modelių ir kainų sąrašas skaitomas iš \code{/v1/models}.

\paragraph{Iškvietimai be būsenos.} Responses API gali laikyti pokalbį tiekėjo pusėje ir leisti
kitai užklausai į jį remtis. Sistema to niekada nenaudoja. \code{gatewayLanguageModel} apgaubia
kiekvieną DI tinklų sąsajos modelį taip, kad kiekvienoje užklausoje parinktis \code{store}
turėtų reikšmę \code{false}: nei DI tinklų sąsaja, nei už jos esantis tiekėjas nesaugo, ko buvo
klausta ir kas atsakyta, o užklausoje visada yra viskas, ką modelis turi perskaityti.
Struktūrizuotam generavimui tai tiesiog viena užklausa ir vienas atsakymas. Pokalbyje tai
reiškia, kad kiekvienas ėjimo žingsnis vėl siunčia visus ankstesnius žingsnius
(\cref{ch:chat}).

\paragraph{Samprotaujantys modeliai.} Kai kurie modeliai prieš atsakydami
\enquote{samprotauja}: jie sukuria vidinius žetonus, kurie apmokestinami kaip išvestis, bet
nėra atsakymo dalis. SDK tokius modelius atpažįsta pagal pačios OpenAI identifikatorius, tačiau
DI tinklų sąsajos identifikatoriai turi diegimo priešdėlį (\code{azure/}), kurio SDK nesitiki.
Todėl \code{isReasoningModel} tą pačią taisyklę taiko pavadinimui po paskutinio pasvirojo
brūkšnio.

\begin{codebox}{Kaip atpažįstamas samprotaujantis modelis ir kas pasikeičia}
DI tinklų sąsajos modelis laikomas samprotaujančiu, kai jo pavadinimas yra \code{gpt-N} arba
\code{gpt-N.M}, kur $N \ge 5$, su bet kokia priesaga -- išskyrus sveikojo skaičiaus versiją,
kurios priesaga prasideda \code{chat}, -- arba kai jis priklauso \code{o} serijai (raidė ir
skaičius). Tokiam modeliui apvalkalas atlieka tris dalykus:
\begin{itemize}
  \item pažymi modelį kaip samprotaujantį (\code{forceReasoning}), kad SDK instrukcijas
    siųstų kaip \enquote{developer} žinutę ir modelio samprotavimą prašytų šifruotu pavidalu
    -- jį vėlesnis pokalbio ėjimo be būsenos žingsnis turi grąžinti atgal;
  \item pašalina iš užklausos \code{temperature} ir \code{topP}, nes samprotaujantis modelis
    nepriima nė vieno iš jų;
  \item palieka \code{store} išjungtą, kaip ir kiekvienam modeliui.
\end{itemize}
Pagal šią taisyklę kiekvienas modelis, kurį nurodo pristatomos instrukcijos
(\cref{sec:models:defaults}), yra samprotaujantis. Todėl instrukcijos versijoje saugoma
temperatūra su pristatomais numatytaisiais nustatymais nesiunčiama; ji įsigalioja tik tada,
kai administratorius instrukciją perkelia į nesamprotaujantį modelį.
\end{codebox}

\paragraph{Atsarginiai modeliai.} Du asistento paslaugos parametrai nurodo modelį, kuris
naudojamas, kai užklausa jo nenurodo: \code{AI_CHAT_MODEL} (numatytasis
\texttt{gateway:azure/gpt-5.5}) generavimui ir \code{AI_EMBED_MODEL} (numatytasis
\texttt{gateway:azure/text-embedding-3-large}) vektoriams. Kiekviena pristatoma instrukcijos
versija nurodo savo modelį, todėl pirmasis -- tik apsauginis tinklas. Antrasis yra kiekvienos
paieškos užklausos vektoriaus modelis; saugomi vektoriai sukuriami modeliu, kurį nurodo paties API to paties
pavadinimo parametras (\code{AI_EMBED_MODEL} faile \file{settings.py}), kurį kviečiantieji
siunčia su kiekvienu paketu, todėl abu parametrai turi nurodyti tą patį modelį
(\cref{sec:models:embed}).


\section{Generavimo maršrutas}\label{sec:models:generate}

\route{POST}{/internal/llm/generate} atlieka vieną generavimą. Tai maršrutas, kuris stovi už
kiekvieno struktūrizuoto atsakymo sistemoje: nuostatų paketo pareigų, verdikto dėl akto,
penkiolikos pareigų išvadų, veiksmų plano. Šiame poskyryje aptariama užklausa, tai, kaip
skaitomas modelis, laiko limitai, schema ir nesėkmės.

\subsection{Ką siunčia API}

Užklausa -- JSON objektas. \code{readGenerateBody} jį patikrina laukas po lauko dar prieš
įtraukiant modelį (\cref{tab:models:request}); nepraėjęs laukas reiškia 400 su atskiru kodu,
ir modelio niekas nepasiekia.

\begin{table}[H]
\caption{Generavimo užklausos laukai}\label{tab:models:request}
\begin{tabularx}{\textwidth}{@{}L{27mm}L{45mm}Y@{}}
\toprule
\hd{Laukas} & \hd{Taisyklė} & \hd{Ką neša} \\
\midrule
\code{prompt} & privalomas, netuščia eilutė (\code{INVALID_PROMPT}) & Šio vieno iškvietimo duomenys: nuostatos, kurias reikia perskaityti, pareigos, kurias reikia įvertinti. Siunčiamas kaip vienintelė naudotojo žinutė. \\
\code{system} & neprivaloma eilutė (\code{INVALID_SYSTEM}) & Instrukcijos: instrukcijos versijos tekstas su užpildytais kintamaisiais. \\
\code{shared} & neprivaloma eilutė (\code{INVALID_SHARED}) & Duomenys, kuriuos nepakeistus neša kiekvienas vieno etapo iškvietimas, pvz., įmonės profilis ir dokumentai ataskaitoje. Dedami po instrukcijų (\cref{sec:models:cache}). \\
\code{model} & neprivaloma eilutė, apkarpoma & Instrukcijos versijos modelio identifikatorius; jei jo nėra -- \code{AI_CHAT_MODEL}. \\
\code{schema} & neprivalomas JSON objektas, ne sąrašas (\code{INVALID_SCHEMA}) & JSON schema, kurią turi atitikti atsakymas. Be jos atsakymas yra paprastas tekstas. \\
\code{schemaName} & neprivaloma eilutė, sutrumpinama iki 64 simbolių; numatytoji \code{result} & Pavadinimas, kuriuo siunčiama schema. \\
\code{temperature} & neprivalomas skaičius nuo 0 iki 2 (\code{INVALID_TEMPERATURE}) & Instrukcijos versijos imties temperatūra; samprotaujančiam modeliui praleidžiama. \\
\code{maxOutputTokens} & neprivalomas sveikasis skaičius nuo 1 iki 128\,000 (\code{INVALID_MAX_TOKENS}) & Instrukcijos versijos išvesties riba, jei ji nustatyta. \\
\bottomrule
\end{tabularx}
\end{table}

\begin{aibox}{Ką modelis gauna struktūrizuoto generavimo metu}
Maršrutas sujungia \code{system} ir \code{shared}, atskirdamas tuščia eilute, ir rezultatą
siunčia kaip užklausos \emph{instrukcijas}; po jų eina \code{prompt} kaip viena naudotojo
žinutė; schema, jei ji yra, siunčiama kaip griežtas atsakymo formatas jos pavadinimu. Tai ir
viskas. Nėra jokios pokalbio istorijos, nėra įrankių ir nesiremiama niekuo iš ankstesnio
iškvietimo: viskas, ką modelis turi žinoti apie teisę, įmonę ar užduotį, yra šiuose trijuose
tekstuose. Ką jie kiekviename kelyje apima, aptariama tolesniuose skyriuose.
\end{aibox}

\subsection{Skaitoma kaip srautas, atsakoma visu atsakymu}

Maršrutas neprašo tiekėjo baigto atsakymo. Jis prašo srauto (SDK \code{streamText}), perskaito
srautą iki galo asistento paslaugos viduje ir tik tada atsako API visu rezultatu viename JSON
objekte:

\begin{codeblock}
{ "model": "gateway:azure/gpt-6.1-sol",
  "text": null,  "object": { ... },      one of the two is null
  "finishReason": "stop",
  "usage": { "inputTokens", "cachedInputTokens", "cacheWriteTokens",
             "outputTokens", "reasoningTokens" },
  "durationMs": 187000 }
\end{codeblock}

Todėl API mato įprastą, tik ilgai trunkančią užklausą. Srautas egzistuoja dėl vienos
priežasties, nurodytos \file{routes/internal.js} antraštiniame komentare: lėtas modelis ilgą
atsakymą rašo daugybę minučių, o Node nutraukia atsakymą, kurio pirmasis baitas atkeliauja
vėliau nei po penkių minučių. Srautas prasideda iš karto, todėl pirmojo baito laukiama
trumpai, kad ir kiek modelis rašytų.

\begin{lessonbox}{Išmatuota 2026 m. lapkričio 10 d.}
Tą dieną DI tinklų sąsajos GPT-6 modeliai rašė apie 70 žetonų per sekundę. Tankus nuostatų
paketas, išsiųstas pareigoms išskirti, užtrukdavo nuo keturių iki šešių minučių; vienas
patikrintas iškvietimas truko 338~sekundes ir parašė 24\,000 išvesties žetonų. Tokios trukmės
iškvietimo neįmanoma atlikti kaip vieno atsakymo, kai pirmajam baitui taikoma penkių minučių
riba, todėl generavimas skaitomas kaip srautas, o toliau aprašytos viršutinės ribos yra
penkiolika minučių.
\end{lessonbox}

Srautas gali sugesti ir jau prasidėjęs -- tiekėjas praneša apie klaidą kaip apie įvykį srauto
viduje. SDK tokį srautą užbaigtų tyliai, ir tai atrodytų kaip trumpesnis atsakymas. Maršrutas
įsimena pirmąją pastebėtą klaidą (\code{onError}) ir, srautui pasibaigus, ją meta:
\enquote{srautas, pranešantis apie klaidą jau prasidėjęs, yra nepavykęs iškvietimas, o ne
trumpesnis atsakymas}, kaip sakoma maršruto antraštiniame komentare. Tai įtvirtina testas
faile \file{tests/internal.test.mjs} (\enquote{joks dalinis atsakymas nepasiekia užduočių
vykdytojo}).

\subsection{Laiko limitai abiejose pusėse}

Abi durų pusės turi suderintas viršutines ribas (\cref{tab:models:timeouts}). Asistento
paslauga vienam generavimui skiria penkiolika minučių; API laukia puse minutės ilgiau, kad
asistento atsisakymas atkeliautų anksčiau, nei API pasiduos, ir priežastis nepasimestų.

\begin{table}[H]
\caption{Modelio iškvietimo laiko limitai}\label{tab:models:timeouts}
\begin{tabularx}{\textwidth}{@{}L{42mm}L{50mm}Y@{}}
\toprule
\hd{Konstanta} & \hd{Kur} & \hd{Reikšmė ir prasmė} \\
\midrule
\code{GENERATE_TIMEOUT_MS} & asistento paslauga, \file{config.js} & 900\,000~ms: vienas generavimas. Jį viršijus, tiekėjo iškvietimas nutraukiamas, o maršrutas atsako 502, nurodydamas laiko limitą. \\
\code{GENERATE_TIMEOUT_S} & API, \file{assistant/gateway.py} & 930~s: kiek \code{generate} laukia asistento paslaugos. \\
\code{EMBED_TIMEOUT_MS} & asistento paslauga, \file{config.js} & 60\,000~ms: vienas vektorizavimo paketas, įskaitant pakartojimus. \\
\code{embed} parametras \code{timeout} & API, \file{assistant/gateway.py} & 120~s: dvigubai daugiau nei asistento riba, kad pakibęs iškvietimas taptų 502, kurį klientas kartoja, o ne veiktų toliau, kai API jau pasidavė. \\
\code{maxRetries} & asistento paslauga, abu maršrutai & 2: SDK nepavykusį tiekėjo iškvietimą pakartoja du kartus (kai biudžetas išnaudotas, klausiama tik vieną kartą). \\
\bottomrule
\end{tabularx}
\end{table}

\subsection{Schema ir jos neatitinkantis atsakymas}

Etapas, kuriam reikia struktūrizuoto atsakymo, siunčia JSON schemą. Schemos priklauso kodui,
kuris atsakymus apdoroja -- kiekvienas konvejerio modulis apibrėžia savąsias, -- o
\file{prompts/schemas.py} jas surašo pagal instrukciją (\code{SCHEMAS}: 21 instrukcija, prie
kiekvienos -- modulis, konstanta ir pavadinimas, kuriuo schema siunčiama), kad bandomasis
paleidimas instrukcijų redaktoriuje naudotų lygiai tą pačią schemą kaip konvejeris. Mažiausioji
iš jų parodo formą:

\begin{codefile}{\file{corpus/citations.py} -- \code{JUDGE_SCHEMA}, siunčiama pavadinimu \code{verdicts}}
JUDGE_SCHEMA = {
    "type": "object", "additionalProperties": False, "required": ["verdicts"],
    "properties": {"verdicts": {"type": "array", "items": {
        "type": "object", "additionalProperties": False, "required": ["ref", "verdict", "note"],
        "properties": {"ref": FILLED,
                       "verdict": {"type": "string", "enum": ["match", "mismatch", "unclear"]},
                       "note": FILLED}}}},
}
\end{codefile}

Kiekviena schema parašyta DI tinklų sąsajos \term{griežtajam režimui}: kiekviena savybė
nurodyta kaip privaloma, \code{additionalProperties} yra \code{false}, o reikšmė, kurios gali
nebūti, deklaruojama kaip tipas, leidžiantis \texttt{null}, o ne kaip neprivalomas raktas. Schema
siunčiama kaip griežtas atsakymo formatas, kuriuo tiekėjo prašoma užtikrinti, kad modelio išvestis jo
laikytųsi: jokio trūkstamo rakto, jokio paties modelio sugalvoto rakto, jokios būsenos už
išvardytųjų ribų.

\paragraph{Privalomas dar nereiškia užpildytas.}\label{par:models:filled} Griežtajame režime
kiekvienas raktas privalomas, tačiau tai užtikrina tik patį raktą: \code{{"type": "string"}}
tenkina ir tuščia eilutė, o masyvą -- ir tuščias sąrašas. Todėl tai, be ko atsakymas negali
apsieiti, pažymima pačioje schemoje (taisyklę nurodo failo \file{prompts/schemas.py}
antraštinis komentaras):

\begin{itemize}
  \item tekstas, kuris privalo turėti turinį, yra \code{FILLED}, eilutė su
    \code{minLength}~1, -- kaip pirmiau parodyti \code{ref} ir \code{note};
  \item sąrašas, kurio dydis žinomas, tai nurodo raktažodžiais \code{minItems} ir
    \code{maxItems}: aštuonios politikos dalys, penki kriterijų lentelės lygiai, bent viena
    nuostata, kuria remiasi pareiga.
\end{itemize}

Laukas, kuriam instrukcija leidžia būti tuščiam, lieka paprasta eilutė: sąlyga, terminas,
dalies numeris. Ribą tarp šių dviejų nubrėžia pačios instrukcijos formuluotė, nes žymė, kurios
instrukcija neparemia, yra blogiau nei jokios žymės: į užpildytu pažymėtą lauką dekoderis
priverčia įrašyti \emph{ką nors}, ir modelis, neturėdamas ko pasakyti, įrašo brūkšnį. DI tinklų
sąsaja užtikrina, kad šių raktažodžių būtų laikomasi, kol modelis rašo, o asistento paslaugos
tikrintuvas juos patikrina dar kartą; abu dalykai išbandyti su visais trimis modeliais prieš
pažymint schemas.

\begin{lessonbox}{Žymė, pakeitusi atsakymus, 2026 m. lapkričio 17 d.}
Kai kurių laukų reikia tik kai kuriuose atsakymuose: rekomendacijos -- kai išvada atvira,
klausimų -- kai akto taikytinumas lieka neaiškus, pastabos -- kai pareiga abejotina. Schema gali
pasakyti ir tai -- dviem elemento formomis po raktažodžiu \code{anyOf}: uždara išvada, kurios
rekomendacija gali būti tuščia, ir atvira išvada, kurios rekomendacija užpildyta, -- ir DI
tinklų sąsaja tokią schemą priima. Ji buvo sukurta, o matavimas nulėmė, kad jos atsisakyta.
Su dviem formomis atitikties spragų analizė sutapo su 64 ir 65 iš 71 vertinimo
rinkinio pažymėto atvejo; su schema, kokia ji buvo iki tol, tą pačią dieną -- su 67 ir 68; su
viena forma ir paprastomis šio poskyrio žymėmis -- su 69 ir 69 (\cref{ch:control}). Du atvejai
pasislinko abiejuose dviejų formų schemos paleidimuose ir nė viename iš kitų keturių, abu
atviros būsenos link: išvada, kurią atvejų žymos laiko atitinkančia, buvo įvertinta kaip
atitinkanti iš dalies, kita -- kaip neįrodyta; kiti skirtumai kiekvieną kartą buvo vis kiti. Niekas dviejose formose teisingo atsakymo nedraudžia; schemos forma yra
dalis to, ką modelis skaito, ir ji pakreipė vertinimą. Trys schemos, turėjusios dvi formas, vėl
turi po vieną, o to, ką jos būtų galėjusios užtikrinti, prašo vien instrukcija. Žymė į schemą
dedama tik ten, kur vertinimas rodo, kad atsakymai su ja nėra prastesni.
\end{lessonbox}

\begin{codebox}{Atsakymas tikrinamas du kartus}
\code{schemaOutput} faile \file{routes/internal.js} perduoda schemą tiekėjui kaip griežtą
atsakymo formatą \emph{ir} iš tos pačios schemos sukuria tikrintuvą, kurį išanalizuotas
atsakymas turi praeiti asistento paslaugos viduje, prieš jį grąžinant. Antraštinis
komentaras nurodo priežastį: tiekėjas be griežtojo režimo gali nukrypti, o API objektą skaito
pagal raktus, todėl pervadintas raktas tyliai reikštų \enquote{išvadų nėra}. Atsakymas,
kurio neįmanoma išanalizuoti kaip JSON arba kuris neatitinka schemos, iškvietimą paverčia
nesėkmingu; kaip rezultatas jis API niekada nepasiekia. Schema, kurios tikrintuvas negali
perskaityti, siunčiama netikrinta, su įspėjimu žurnale, ir tada už jos stovi tik tiekėjo
griežtasis režimas.
\end{codebox}

Nepavykusi patikra neturi tiekėjo HTTP būsenos, todėl maršrutas į ją atsako
502~\code{GENERATION_FAILED}, kurį API klientas laiko laikinu sutrikimu ir klausia dar kartą
(\cref{sec:models:client}). Schema, kurią atmeta pati DI tinklų sąsaja -- t.~y. neatitinkanti
griežtosios formos, -- grįžta kaip pačios sąsajos 4xx ir nekartojama: klausti dar kartą
nepadės.

\subsection{Kelio kortelė}

Šiuo keliu eina kiekvienas struktūrizuotas generavimas sistemoje. Vėlesnių skyrių kortelės
nurodo, kas būdinga kiekvienai instrukcijai; ši nurodo, kas bendra visoms.

\begin{pathcard}{Struktūrizuotas generavimas}{run_prompt}
  \pcrow{Trigger}{Užduočių vykdytojo užduoties etapas arba maršrutas, kurio laukia žmogus,
    kviečia API funkciją \code{run_prompt} (\file{assistant/llm.py}) su instrukcijos
    identifikatoriumi (slug) ir duomenimis. Instrukcijų redaktoriaus bandomasis paleidimas
    kviečia \code{gateway.generate} tiesiogiai su projekto tekstu.}
  \pcrow{Input}{Trys tekstai: instrukcijos (aktyvi instrukcijos versija su užpildytais
    \code{{PLACEHOLDER}} kintamaisiais); nebūtinai -- bendrasis blokas, kurį neša visi etapo
    iškvietimai; šio iškvietimo duomenys. Be to, JSON schema ir jos pavadinimas, modelio
    identifikatorius, instrukcijos versijos temperatūra ir išvesties riba. Turinys gali būti
    iki 16~MB; paketų dydžius ir dydžio ribas nustato kiekvienas kelias.}
  \pcrow{Prompt and model}{Instrukcija, nurodyta identifikatoriumi; jos aktyvios versijos
    modelis, nebent kviečiantysis perduoda kitą. Užduoties etapui atliekamas iškvietimas yra
    \emph{kantrus}: jis išlaukia dažnio limito. Samprotaujančiam modeliui temperatūra
    nesiunčiama.}
  \pcrow{Output}{Su schema -- \code{object}: išanalizuotas JSON. Be jos -- \code{text}.
    Visada \code{model}, \code{finishReason}, \code{durationMs} ir penki \code{usage} žetonų
    skaičiai.}
  \pcrow{Checked by code}{Užklausa, laukas po lauko. Schemą užtikrina tiekėjo griežtasis
    režimas, o asistento paslauga ją patikrina dar kartą. Klaida srauto viduje iškvietimą
    paverčia nesėkmingu. Visa, kas susiję su atsakymo \emph{turiniu}, tikrina kviečiantysis
    kelias.}
  \pcrow{Stored}{Viena eilutė lentelėje \code{llm_calls}: rūšis \code{job_step}
    (\code{prompt_test} bandomajam paleidimui), paskirtis, modelis, penki žetonų skaičiai,
    kaina, trukmė, užduoties etapas, apmokestintas naudotojas ir instrukcijos versija. Patį
    atsakymą saugo kviečiantysis. Šiame sluoksnyje sistema atsakymų podėlio nelaiko;
    instrukcijas laiko tiekėjo podėlis (\cref{sec:models:cache}).}
  \pcrow{On failure}{SDK nepavykusį tiekėjo iškvietimą pakartoja du kartus, išnaudoto
    biudžeto -- niekada. API klientas laikiną sutrikimą pakartoja vieną kartą; užduoties
    etapas dažnio limito laukia iki penkiolikos minučių; nutrūkęs ryšys bandomas penkis
    kartus. Tada kviečiantįjį pasiekia \code{GatewayError} (\cref{tab:models:errors}).
    Nepavykęs iškvietimas sąnaudų eilutės nepalieka.}
  \pcrow{Cost}{Apskaičiuojama dalis po dalies pagal DI tinklų sąsajos kainoraštį ir
    priskiriama naudotojui, kurį nurodo kviečiantysis, arba niekam
    (\cref{tab:models:charged}). Išmatuota 2026 m. lapkričio 10--11 d. su pavyzdine įmone:
    atitikties spragų analizės iškvietimas su 59\,000 įvesties žetonų kainavo 0,18~\$, kai
    pildė tiekėjo podėlį, ir apie 0,045~\$, kai skaitė iš jo.}
  \pcrow{Human control}{Šiame sluoksnyje -- jokios. Administratorius parenka kiekvienos
    instrukcijos tekstą ir modelį (\cref{ch:prompts}) bei kiekvienos paskyros dienos
    biudžetą; ką žmogus peržiūri atsakyme, nurodyta kiekvieno kelio kortelėje.}
  \pcrow{Measured by}{Testai, ne vertinimai: asistento paslaugos
    \file{tests/internal.test.mjs}, \file{tests/provider.test.mjs} ir
    \file{tests/usage.test.mjs} (tinklo formatas, schemos patikra, nesėkmės, sąnaudos); API
    \file{tests/test_prompts.py} (kainos skaičiavimas, pakartojimai, kantrybė, kainoraštis).
    Niekas šiame sluoksnyje nematuoja instrukcijos.}
\end{pathcard}


\section{Vektorizavimo maršrutas}\label{sec:models:embed}

\term{Vektorinis atvaizdas}, trumpiau -- vektorius (angl. \en{embedding}), yra iš teksto
apskaičiuota skaičių seka, tokia, kad panašios prasmės tekstai gautų arti vienas kito esančius
vektorius. Sistema saugo po
vektorių kiekvienam teisės aktų bazės fragmentui, kiekvienai pareigai ir kiekvienam įmonės
dokumentų fragmentui, o kiekvienai paieškos užklausai jį apskaičiuoja;
\cref{ch:corpus,ch:company} aprašo, kas su jais daroma. Visi jie gaunami iš
\route{POST}{/internal/llm/embed}.

\paragraph{Užklausa.} Sąrašas \code{texts} iš 1--256 eilučių (\code{MAX_EMBED_TEXTS}; kitaip
400~\code{INVALID_TEXTS} arba \code{TOO_MANY_TEXTS}) ir nebūtinai \code{model}; jei jo nėra,
maršrutas naudoja \code{AI_EMBED_MODEL}. Dėl tuščios eilutės kai kurie tiekėjai atmestų visą
paketą, todėl ji pakeičiama vienu tarpu. SDK \code{embedMany} paketą siunčia daugiausia
keturiais lygiagrečiais iškvietimais ir su dviem pakartojimais, neviršydama ribos
\code{EMBED_TIMEOUT_MS}.

\paragraph{Modelis ir jo plotis.} Vektorizavimo modelis -- DI tinklų sąsajos
\model{text-embedding-3-large}. Jo vektorių prašoma \textbf{1536 matmenų}, o ne viso modelio
pločio, nes toks yra kiekvieno vektorių stulpelio duomenų bazėje plotis
(\code{EMBED_DIMENSIONS} failuose \file{corpus/models.py} ir \file{companies/models.py},
pakartota faile \file{routes/internal.js}). Maršrutas parinktį \code{dimensions} perduoda tik
\code{text-embedding-3} šeimos modeliams, t.~y. tiems, kurie ją priima. Kaip antroji
apsauga, \code{fitDimensions} bet kurį ilgesnį vektorių sutrumpina iki pirmųjų 1536 reikšmių ir
iš naujo normalizuoja -- tą patį daro ir parinktis, -- kad tilptų ir parinkties nepaisęs
tiekėjas. Trumpesnis vektorius perduodamas toliau, ir jį atmeta duomenų bazės stulpelis.

Kadangi kiekvienas saugomas vektorius, kad būtų palyginamas, turi būti gautas iš to paties
modelio, pakeitus \code{AI_EMBED_MODEL} viską tenka vektorizuoti iš naujo (užduotis
\code{corpus.reembed}; \cref{ch:corpus}).

\paragraph{Atsakymas.} \code{embeddings} (po vieną vektorių kiekvienam tekstui, ta pačia
tvarka, kuria buvo išsiųsti), \code{model}, \code{dimensions}, \code{usage.tokens} ir
\code{durationMs}.

\paragraph{Paketai.} Maršruto riba -- 256 tekstai; kviečiantieji lieka gerokai žemiau jos
(\cref{tab:models:embedcallers}). Kiekvienas paketas -- vienas iškvietimas ir viena sąnaudų
eilutė.

\begin{table}[H]
\caption{Kas prašo vektorių ir kokiais paketais}\label{tab:models:embedcallers}
\begin{tabularx}{\textwidth}{@{}L{43mm}YL{40mm}@{}}
\toprule
\hd{Kviečiantysis} & \hd{Vektorizuojamas tekstas} & \hd{Paketas; užrašoma paskirtis} \\
\midrule
\file{corpus/ingest.py}, \code{embed_chunks} & Akto fragmentas, pradedamas trumpuoju akto pavadinimu, nuostatos keliu ir jos antrašte & 64 (\code{EMBED_BATCH}); \code{corpus.embed} \\
\file{corpus/obligations.py}, \code{embed_obligations} & Pareigos pavadinimas ir tekstas, o jei yra -- ir angliškas pavadinimas bei tekstas & 32 pareigos, daugiausia 64 tekstai; \code{obligations.embed} \\
\file{corpus/obligations.py}, \code{embed_english} & Pareigų, kurioms trūksta šio vektoriaus, angliškas pavadinimas ir tekstas & 64; \code{obligations.embed_en} \\
\file{companies/indexing.py}, \code{index_document} & Įmonės dokumento langas, pradedamas dokumento pavadinimu ir skirsnio antrašte & 64; \code{documents.embed} \\
\file{corpus/search.py}, \file{companies/search.py} & Paieškos užklausa & 1; neužrašoma \\
\bottomrule
\end{tabularx}
\end{table}

\begin{pathcard}{Vektorizavimo iškvietimas}{gateway.embed}
  \pcrow{Trigger}{Akto įkėlimas arba pakartotinis vektorizavimas (užduotys
    \code{corpus.ingest}, \code{corpus.reembed}); pareigų išskyrimas arba administratoriaus
    atliktas pareigos taisymas; įkelto dokumento indeksavimas (užduotis
    \code{documents.index}); ataskaita, aptikusi dar neindeksuotą dokumentą arba pareigas be
    angliško vektoriaus; kiekviena paieška teisės aktų bazėje ar įmonės dokumentuose.}
  \pcrow{Input}{Tekstų sąrašas, daugiausia 256 per iškvietimą (\code{MAX_EMBED_TEXTS});
    kviečiantieji siunčia daugiausia 64 arba vieną užklausą (\cref{tab:models:embedcallers}).
    Kiekvienas tekstas -- ištrauka kartu su tuo, kas nurodo jos vietą: aktu ir nuostata arba
    dokumentu ir skirsniu.}
  \pcrow{Prompt and model}{Instrukcijos nėra. \texttt{gateway:azure/text-embedding-3-large}
    (\code{AI_EMBED_MODEL}), prašoma 1536 matmenų.}
  \pcrow{Output}{\code{embeddings} -- po vieną 1536 skaičių vektorių kiekvienam tekstui ta
    tvarka, kuria išsiųsta; \code{model}; \code{dimensions}; \code{usage.tokens};
    \code{durationMs}.}
  \pcrow{Checked by code}{Tekstai patikrinami, o tušti pakeičiami tarpu. Ilgesnis vektorius
    sutrumpinamas iki 1536 reikšmių ir iš naujo normalizuojamas. Teisės aktų įkėlimas atmeta
    atsakymą, kurio vektorių skaičius skiriasi nuo tekstų skaičiaus. Kito pločio vektorių
    atmeta duomenų bazės stulpelis.}
  \pcrow{Stored}{Vektorių stulpeliai \code{legal_chunks.embedding},
    \code{obligations.embedding} ir \code{obligations.embedding_en},
    \code{company_document_chunks.embedding}. Užklausos vektorius panaudojamas vieną kartą ir
    nesaugomas. Po vieną \code{embed} rūšies \code{llm_calls} eilutę kiekvienam paketui,
    išskyrus užklausą.}
  \pcrow{On failure}{SDK -- du pakartojimai; paketas, užtrukęs ilgiau nei 60~s, nutraukiamas
    ir į jį atsakoma 502. API klientas laikiną sutrikimą kartoja tris kartus (po 2, 6 ir
    18~s). Paieška, kurios užklausos nepavyksta vektorizuoti, vykdoma vien pagal raktažodžius.
    Indeksavimo užduotis bandoma tris kartus.}
  \pcrow{Cost}{0,13~\$ už milijoną žetonų (DI tinklų sąsajos kainoraštis, 2026 m. lapkričio
    10 d.). Įkelto dokumento vektorių kaina priskiriama įkėlusiajam; teisės aktų bazės darbas --
    niekam; užklausa neužrašoma. Išmatuota: 24 naujų aktų įkėlimas duomenų bazės kopijoje
    2026 m. lapkričio 12 d. kainavo 0,08~\$; visko pakartotinis vektorizavimas 2026 m.
    rugsėjo 28 d. įvertintas maždaug 0,30~\$.}
  \pcrow{Human control}{Jokios.}
  \pcrow{Measured by}{\file{tests/internal.test.mjs} ir \file{tests/embed-timeout.test.mjs}
    (maršrutas, 1536 matmenys, viršutinė riba); 2026 m. rugsėjo 28 d. dviejų vektorizavimo
    modelių palyginimas (\cref{tab:models:embedcompare}).}
\end{pathcard}


\section{API klientas}\label{sec:models:client}

API pusėje su asistento paslauga kalbasi vienas modulis -- \file{assistant/gateway.py}, o virš
jo yra viena funkcija, kurią kviečia konvejeriai, -- \code{run_prompt} faile
\file{assistant/llm.py}. Modulis vadinamas \enquote{gateway}, nes tai API vartai į modelius;
jo nereikia painioti su fakulteto DI tinklų sąsaja, į kurią jis niekada nesikreipia.

\subsection{Du iškvietimai}

\code{gateway.generate} iš savo argumentų sudaro užklausą (\cref{tab:models:request}),
praleidžia tuos, kurie yra \code{None}, išsiunčia ją su bendrąja paslaptimi ir grąžina
\code{GenerateResult}: tekstą arba objektą, modelį, pabaigos priežastį, penkis žetonų skaičius
ir trukmę. \code{gateway.embed} tą patį daro su tekstų sąrašu ir grąžina \code{EmbedResult}.
Trečioji funkcija, \code{list_models}, perskaito modelių sąrašą ir, kai asistento paslauga
neveikia, grąžina tuščią, kad niekas, kam reikia tik kainų, dėl to nesugestų.

Abu iškvietimai eina per \code{_post} -- ciklą, kuris siunčia užklausą ir sprendžia, ką daryti
su nesėkme (\cref{sec:models:errors}). Galutinė nesėkmė iškeliama kaip \code{GatewayError} su
trimis dalykais, pagal kuriuos kviečiantysis gali veikti: HTTP būsena (0 tinklo sutrikimui),
mašininis kodas ir asistento paslaugos pranešimas, kuriame yra paties tiekėjo žodžiai. Tinklo
klaidos tekste minimi vidiniai serveriai ir prievadai; jis keliauja į žurnalą, o klaida, kurią
naudotojas gali perskaityti ataskaitoje ar užduotyje, žodžiais pasako, kas atsitiko
(\enquote{DI paslauga neatsakė per 930~s}).

\subsection{run\_prompt: saugoma instrukcija virš duomenų}

\code{run_prompt}, kaip sakoma jos antraštiniame komentare, yra \enquote{vienintelis būdas,
kuriuo konvejerio etapas kalbasi su modeliu}. Jos signatūra rodo, ką sprendžia kviečiantysis
ir ko jis nesprendžia:

\begin{codeblock}
run_prompt(slug, data, *, shared=None, schema=None, schema_name="result",
           variables=None, job_step_id=None, user_id=None, purpose=None, model=None)
\end{codeblock}

Ji atlieka keturis veiksmus iš eilės.

\begin{enumerate}
  \item \textbf{Randa instrukcijos versiją.} \code{active_version(slug)}
    (\file{prompts/store.py}) grąžina versiją, kurią administratorius tai instrukcijai padarė
    aktyvia. Jei tokios nėra, ji iškelia \code{PromptNotConfigured}, ir etapas nepavyksta su
    tuo pranešimu: etapas niekada nevykdomas pagal instrukcijas, kurių niekas negali matyti.
  \item \textbf{Užpildo kintamuosius.} \code{render} pakeičia kiekvieną \code{{NAME}} formos
    kintamąjį (didžiosios raidės, skaitmenys, pabraukimo brūkšniai), kuriam kviečiantysis
    pateikė reikšmę; ataskaita perduoda \code{LANGUAGE}, \code{TODAY} ir
    \code{COMPANY_NAME}. Kintamasis be reikšmės lieka tekste toks, koks parašytas, todėl
    redaktoriuje su klaida įrašytas pavadinimas matomas instrukcijose, o ne dingsta.
  \item \textbf{Iškviečia modelį.} Užpildytas tekstas yra instrukcijos; \code{data} -- naudotojo
    žinutė; \code{shared}, schema ir jos pavadinimas perduodami toliau; modelis, temperatūra ir
    išvesties riba imami iš instrukcijos versijos (kviečiantysis gali pakeisti modelį).
    Iškvietimas yra \emph{kantrus} būtent tada, kai nurodytas \code{job_step_id}.
  \item \textbf{Užrašo iškvietimą} su \code{record_call} (\cref{sec:models:usage}), įskaitant
    instrukcijos versijos identifikatorių, \enquote{kad ataskaitą visada būtų galima atsekti iki
    tikslių instrukcijų, pagal kurias ji parengta}.
\end{enumerate}

Ji grąžina \code{PromptRun}: objektą arba tekstą, neapdorotą rezultatą ir panaudotą versiją.
Kviečiantysis pasirenka duomenis, schemą, kam priskiriama kaina ir prie kurio užduoties etapo
iškvietimas registruojamas; instrukcijų teksto ir modelio kviečiančiojo kode niekada nėra. Būtent
tai leidžia administratoriui pakeisti bet kurį iš jų be diegimo (\cref{ch:prompts}).

\subsection{Bendrieji duomenys ir lygiagretūs iškvietimai}

Dvi pagalbinės funkcijos lemia, kaip etapas atlieka daug iškvietimų.

\code{parallel_map(fn, items, workers=4)} faile \file{common/parallel.py} vykdo etapo
iškvietimus po kelis vienu metu gijose ir rezultatus grąžina elementų tvarka. Beveik visą laiką
etapas laukia modelio, todėl keturi iškvietimai vienu metu ataskaitą sutrumpina nuo kelių
dešimčių minučių iki kelių. Kiekviena gija, baigusi darbą, uždaro savo duomenų bazės ryšį. Jei
iškvietimas nepavyksta, kitiems vis tiek leidžiama baigtis, o pirmoji išimtis iškeliama, kai
visi užsibaigia.

\code{shared=} pažymi tą įvesties dalį, kuri visuose etapo iškvietimuose ta pati: ataskaitos
atitikties spragų analizėje tai įmonės profilis ir jos dokumentai. \code{run_prompt} ją
perduoda atskirai nuo duomenų, o asistento paslauga ją deda po instrukcijų, kur ją išsaugo
tiekėjo podėlis. Kad tai duotų naudos, pirmasis iškvietimas turi būti baigtas prieš pradedant
kitus -- kartu pradėti iškvietimai kiekvienas mokėtų už bloko įrašymą. \code{shared_first}
faile \file{analysis/pipeline.py} būtent tai ir daro: pirmąjį elementą vykdo vieną, o likusius
perduoda \code{parallel_map}. \Cref{sec:models:cache} pateikia tai pagrindžiantį matavimą.

\subsection{Pakartojimai, kantrybė ir klaidos, kurias mato kviečiantysis}\label{sec:models:errors}

Nesėkmė apdorojama trimis lygmenimis, kurių kiekvieno horizontas kitoks. SDK asistento
paslaugoje nepavykusį tiekėjo iškvietimą per kelias sekundes pakartoja du kartus. Tada API
klientas renkasi tarp trumpo pakartojimo ir ilgo laukimo. Kas lieka, pasiekia kviečiantįjį
kodą kaip \code{GatewayError}.

Klientas skiria dvi kviečiančiųjų rūšis. Iškvietimas, kurio laukia žmogus -- nuorodų patikra
iš teisės aktų bazės puslapio, bandomasis paleidimas, -- turi nepavykti pakankamai greitai,
kad tą žmogų būtų galima informuoti. Užduoties etapui atliekamas iškvietimas yra foninis
darbas, kurio niekas nestebi: jam dažnio limitas yra \enquote{visai ne nesėkmė}, nes tiekėjas
skaičiuoja žetonus per minutę, ir limitas praeina savaime. \Cref{tab:models:retries} pateikia
\code{_post} taisykles.

\begin{table}[H]
\caption{Ką API klientas daro su nepavykusiu iškvietimu. $n$ -- jau atliktų pakartojimų
skaičius; \code{generate} pagal numatymą kviečiama su vienu pakartojimu, \code{embed} -- su
trimis, instrukcijų redaktoriaus bandomasis paleidimas -- be
pakartojimų.}\label{tab:models:retries}
\begin{tabularx}{\textwidth}{@{}L{44mm}Y@{}}
\toprule
\hd{Nesėkmė} & \hd{Ką daro klientas} \\
\midrule
Dažnio limitas (429), kantrus iškvietimas & Laukia ir klausia dar kartą, neišnaudodamas pakartojimo: tiek, kiek prašė tiekėjas (\code{retryAfterS}), kitaip \code{RATE_LIMIT_WAIT_S} = 20~s, paskui 40, paskui daugiausia 60; kiekvienas laukimas atsitiktinai ketvirčiu ilgesnis ar trumpesnis, kad vieno etapo iškvietimai negrįžtų visi tą pačią sekundę. Iš viso išlaukus \code{RATE_LIMIT_PATIENCE_S} = 900~s, limitas traktuojamas kaip kitoje eilutėje. \\
Dažnio limitas (429), kiti iškvietimai; 408, 409, 425, 500, 502, 503, 504 & Kartojama iki kviečiančiojo nurodyto pakartojimų skaičiaus, po $\min(2\cdot 3^{n},\,30)$ sekundžių: 2, 6, 18. \\
Nepavyko prisijungti (\code{GATEWAY_UNREACHABLE}) & Kartojama ne mažiau kaip \code{CONNECTION_RETRIES} = 5 kartus, jei kviečiantysis leidžia bent vieną pakartojimą (po 2, 6, 18, 30 ir 30~s): asistento paslauga buvo paleidžiama iš naujo, o pats iškvietimas buvo tvarkingas. \\
Laiku negautas atsakymas (\code{GATEWAY_TIMEOUT}) & Kartojama iki kviečiančiojo nurodyto pakartojimų skaičiaus. \\
Bet kuri kita 4xx, įskaitant išnaudotą biudžetą (402) & Iškeliama iš karto. Klausti dar kartą nepadės. \\
\bottomrule
\end{tabularx}
\end{table}

\Cref{tab:models:errors} išvardija viską, ką kviečiantysis gali gauti. Asistento paslauga šiuos
atsakymus sudaro funkcijoje \code{providerFailure} (\file{routes/internal.js}) pagal
klasifikaciją faile \file{middleware/errors.js}: išnaudotas biudžetas tampa 402, dažnio limitas
-- 429 su tiekėjo nurodytu laukimu, o visa kita išlaiko paties tiekėjo būseną arba tampa 502,
jei jos nebuvo, -- kad klientas galėtų atskirti kartotiną nesėkmę nuo nekartotinos.

\begin{xltabular}{\textwidth}{@{}L{41mm}L{50mm}Y@{}}
\caption{Modelio iškvietimo klaidos, kaip jas gauna API}\label{tab:models:errors}\\
\toprule
\hd{Būsena ir kodas} & \hd{Kas atsitiko} & \hd{Kas toliau} \\
\midrule
\endfirsthead
\toprule
\hd{Būsena ir kodas} & \hd{Kas atsitiko} & \hd{Kas toliau} \\
\midrule
\endhead
\bottomrule
\endlastfoot
402 \code{AI_BUDGET_EXHAUSTED} & DI tinklų sąsaja atmetė iškvietimą, nes jos rakto biudžetas išnaudotas (\cref{sec:models:budgets}); taip pat tiekėjo žodžiais išreikštas kvotos atsisakymas. & SDK klausia vieną kartą, klientas niekada nekartoja. Ataskaita nepavyksta su pranešimu, prasidedančiu \code{AI_BUDGET_EXHAUSTED}, kurį konsolė parodo žodžiais, ir išsaugo jau apmokėtą darbą. \\
429 \code{GENERATION_FAILED} arba \code{EMBEDDING_FAILED} & Tiekėjo dažnio limitas truko ilgiau nei du SDK pakartojimai arba atėjo srauto viduje. \code{retryAfterS} neša laukimą, jei tiekėjas jį nurodė. & Užduoties etapas jo laukia iki penkiolikos minučių; iškvietimas, kurio laukia žmogus, trumpai kartojamas (\cref{tab:models:retries}). \\
400 \code{UNKNOWN_MODEL} & Modelio identifikatorius nenurodo jokio modelio, nurodo nežinomą tiekėją arba tiekėją be prisijungimo duomenų. & Iškeliama iš karto. Etapas nepavyksta su pranešimu. \\
400, kodas, prasidedantis \code{INVALID_} & Užklausos laukas nepraėjo patikros (\cref{tab:models:request}). & Iškeliama iš karto; kviečiančiojo kodo defektas. \\
4xx \code{GENERATION_FAILED} & Tiekėjas atmetė pačią užklausą: negriežta schema, modelis, kurio raktas negali naudoti, įvestis, ilgesnė, nei modelis priima. & Iškeliama iš karto, pranešime -- tiekėjo žodžiai. \\
5xx \code{GENERATION_FAILED} & Paties tiekėjo klaida po SDK pakartojimų. & 500, 502, 503 ar 504 klientas pakartoja, tada iškelia. \\
502 \code{GENERATION_FAILED} arba \code{EMBEDDING_FAILED} & Tiekėjo būsenos nėra: atsakymas nebuvo galiojantis JSON arba neatitiko schemos, iškvietimas pasiekė savo laiko ribą arba sutriko tinklas iki tiekėjo. & Klientas pakartoja, tada iškelia. \\
0 \code{GATEWAY_TIMEOUT} & Asistento paslauga neatsakė per kliento laiko limitą. & Pakartojama, tada iškeliama. \\
0 \code{GATEWAY_UNREACHABLE} & Nėra ryšio su asistento paslauga. & Pakartojama penkis kartus (kviečiančiajam, kuris neleidžia pakartojimų, -- nė karto), tada iškeliama. \\
403 \code{INTERNAL_ONLY} & Bendrosios paslapties nėra arba ji neteisinga. & Iškeliama iš karto; konfigūracijos defektas. \\
\code{PromptNotConfigured} & Ne iškvietimo klaida: instrukcija neturi aktyvios versijos. & \code{run_prompt} iškelia ją prieš bet kokį iškvietimą. \\
\end{xltabular}

Ką kviečiantysis daro su iškelta klaida, priklauso nuo to, ar jis gali apsieiti be modelio.
Ataskaitos etapai negali: išimtis užbaigia užduotį, o ataskaita užrašo pranešimą
(\code{failure_message} faile \file{analysis/pipeline.py}; \cref{ch:report}). Keli mažesni
skaitytojai gali, ir kiekvienas turi nustatytą atsarginį sprendimą:

\begin{itemize}
  \item nuorodų patikra vertina pagal žodžių sutapimą, o ne modeliu
    (\file{corpus/citations.py}) -- išskyrus tai, kad ataskaitos viduje 402 iškeliama kaip bet
    kuri kita etapo nesėkmė, kad be biudžeto parengta ataskaita nebūtų palaikyta išsamia;
  \item standarto struktūros skaitymas grįžta vien prie struktūros taisyklių, o kai šios nieko
    neranda -- prie standarto skaidymo pagal puslapius (\code{read_standard} faile
    \file{corpus/outline.py});
  \item akto apimties ir jo sankcijų nuostatų skaitymas užrašo nesėkmę žurnale ir palieka aktą
    neperskaitytą, kad vėlesnis paleidimas klaustų dar kartą (\file{corpus/scope.py},
    \file{corpus/sanctions.py});
  \item paieška, kurios užklausos nepavyksta vektorizuoti, vykdoma vien pagal raktažodžius
    (\file{corpus/search.py}, \file{companies/search.py});
  \item instrukcijų redaktoriaus bandomasis paleidimas atsako 502~\code{MODEL_ERROR} su
    priežastimi.
\end{itemize}

Kiekvieno kelio skyriai nurodo, ką atsarginis sprendimas reiškia rezultatui.


\section{Sąnaudos: kas užrašoma apie kiekvieną iškvietimą}\label{sec:models:usage}

Kiekvienas modelio iškvietimas palieka vieną eilutę lentelėje \code{llm_calls}. Eilutę įrašo
\code{record_call} faile \file{assistant/usage.py}, kurią naudoja keturių rūšių
kviečiantieji: \code{run_prompt} -- kiekvienam struktūrizuotam generavimui, instrukcijų
redaktoriaus bandomasis paleidimas (\code{test_run} faile \file{prompts/api/views.py}),
vektorių kviečiantieji -- kiekvienam paketui, ir vidinis maršrutas
\route{POST}{/internal/assistant/turn-log}, per kurį asistento paslauga užrašo kiekvieną
pokalbio ėjimą ir kiekvieną savo parašytą pokalbio pavadinimą.

\subsection{Penki skaičiai}

Iškvietimo kaina priklauso ne vien nuo to, kiek žetonų įėjo, nes tiekėjas tris įvesties rūšis
apmokestina trimis kainomis (\cref{sec:models:cache}). Todėl iškvietimas skaičiuojamas penkiais
skaičiais, kuriuos iš SDK pateikiamų duomenų paima \code{usageOf} faile \file{llm/usage.js}:

\begin{enumerate}
  \item \code{inputTokens} -- visa įvestis;
  \item \code{cachedInputTokens} -- įvesties dalis, kurią tiekėjas nuskaitė iš savo
    instrukcijų podėlio;
  \item \code{cacheWriteTokens} -- įvesties dalis, kurią tiekėjas įrašė į tą podėlį;
  \item \code{outputTokens} -- visa išvestis;
  \item \code{reasoningTokens} -- išvesties dalis, kuri buvo modelio samprotavimas.
\end{enumerate}

Tiekėjas, kuris apie podėlį nepraneša, antrąjį ir trečiąjį skaičius palieka nuliais. Pokalbio
ėjimui, kurį sudaro keli modelio žingsniai, penki skaičiai yra žingsnių sumos
(\code{addUsage}).

\subsection{Iškvietimo kaina}

\code{cost_of} skaičius paverčia JAV doleriais. Tegu $I$ -- įvesties žetonai, iš kurių $c$
nuskaityti iš podėlio, o $w$ į jį įrašyti, $O$ -- išvesties žetonai, ir tegu modelis turi
keturias kainas už žetoną: $p$ -- įvesčiai, $p_r$ -- iš podėlio nuskaitytai įvesčiai, $p_w$ --
į podėlį įrašytai įvesčiai, $p_o$ -- išvesčiai. Tada
\[
  \mathrm{kaina} \;=\; p\,(I - c - w) \;+\; p_r\,c \;+\; p_w\,w \;+\; p_o\,O ,
\]
suapvalinta iki milijonosios dolerio dalies. Trys detalės užtikrina, kad skaičius būtų sąžiningas. Abi dalys
apribojamos taip, kad niekada negalėtų sudaryti daugiau nei visa įvestis ($c \le I$,
$w \le I - c$). Modeliui, kuriam kainoraštis nenurodo podėlio kainos, ta dalis apmokestinama
paprasta įvesties kaina. O modelis visai be kainos kainuoja 0, ir tai užrašoma žurnale vieną
kartą kiekvienam modeliui -- sąnaudų eilutė vis tiek įrašoma su savo žetonais.

Pavyzdys su \model{gpt-6.1-sol} kainomis (tą pačią aritmetiką įtvirtina testas faile
\file{tests/test_prompts.py}): iškvietimas su 60\,000 įvesties žetonų, iš kurių 55\,000
nuskaityti iš podėlio, o 4\,000 į jį įrašyti, ir 3\,000 išvesties žetonų kainuoja
$1\,000 \times 2 + 55\,000 \times 0{,}10 + 4\,000 \times 2{,}50 + 3\,000 \times 10$ milijonųjų
dolerio dalių, t.~y. 0,0475~\$. Įkainotas vien pagal žetonų skaičius, jis būtų užrašytas kaip
0,15~\$.

\subsection{Kainoraštis}

Kainos į kodą neįrašomos kaip tiesa; jos skaitomos iš DI tinklų sąsajos. \code{price_table}
paprašo asistento paslaugos modelių sąrašo, kuriame prie kiekvieno modelio nurodyta įvesties ir
išvesties žetono kaina, o kur DI tinklų sąsaja jas nurodo, -- ir dvi podėlio kainos. Lentelė
laikoma atmintyje valandą (\code{PRICE_TTL_S} = 3600; pati asistento paslauga DI tinklų sąsajos
sąrašą laiko podėlyje dešimt minučių). Už jos stovi trys sluoksniai, kad DI tinklų sąsajos
trikdis nei pamestų sąnaudų eilutės, nei nunulintų biudžeto:

\begin{enumerate}
  \item kiekvienas sėkmingas gavimas taip pat išsaugomas lentelėje \code{model_prices} (po
    eilutę kiekvienam modeliui: keturios kainos ir laikas);
  \item nepavykęs gavimas grįžta prie tos išsaugotos lentelės, kuri tada valandą tarnauja kaip
    šviežia; tik kai ir ten nieko nėra, gauti bandoma vėl po penkių minučių
    (\code{PRICE_RETRY_S} = 300);
  \item po abiem slypi \code{KNOWN_PRICES} -- kainos, žinomos rašant kodą
    (\cref{tab:models:prices}), skirtos diegimui, kuris niekada nepasiekė DI tinklų sąsajos,
    ir modeliui, kurio gautas sąrašas nenurodo.
\end{enumerate}

\begin{table}[H]
\caption{Kainos, kurias kodas pateikia kaip paskutinę atsarginę priemonę (\code{KNOWN_PRICES}
faile \file{assistant/usage.py}), JAV doleriais už milijoną žetonų. Tai DI tinklų sąsajos
kainos 2026 m. rugsėjo 27 d.; GPT-6 modelių ir visos podėlio kainos -- 2026 m. lapkričio
10 d. Brūkšnys reiškia, kad sąraše tokios kainos nėra ir taikoma įvesties
kaina.}\label{tab:models:prices}
\begin{tabularx}{\textwidth}{@{}Yrrrr@{}}
\toprule
\hd{Modelis} & \hd{Įvestis} & \hd{Išvestis} & \hd{Nuskaityta iš podėlio} & \hd{Įrašyta į podėlį} \\
\midrule
\model{azure/gpt-6-astra} & 10 & 50 & 1,00 & 12,50 \\
\model{azure/gpt-6.1-sol} & 2 & 10 & 0,10 & 2,50 \\
\model{azure/gpt-6-sol} & 2 & 10 & 0,20 & 2,50 \\
\model{azure/gpt-5.5} & 5 & 30 & 0,50 & --- \\
\model{azure/gpt-5.6-sol} & 4 & 20 & --- & --- \\
\model{azure/gpt-5.6-terra} & 2 & 12 & 0,20 & 2,50 \\
\model{azure/gpt-5.6-luna} & 0,20 & 1,20 & 0,02 & 0,25 \\
\model{azure/text-embedding-3-large} & 0,13 & 0 & --- & --- \\
\model{azure/text-embedding-3-small} & 0,02 & 0 & --- & --- \\
\bottomrule
\end{tabularx}
\end{table}

Įvertis, rodomas prieš pradedant ataskaitą (\file{analysis/estimate.py}; \cref{ch:report}),
naudoja tą patį \code{cost_of} ir tą pačią lentelę, kiekvieno etapo aktyvios instrukcijos
versijos modeliui. Ataskaitos etapų įvestį jis įkainoja kaip įrašytą į podėlį -- tiek tiekėjas
ima už įvestį, kurią mato pirmą kartą, -- o atitikties spragų analizę kaip vieną iškvietimą,
kuris įrašo bendrąjį bloką, ir likusius, kurie jį nuskaito; vienkartinis dar neperskaityto akto
pareigų išskyrimas įkainojamas paprasta įvesties kaina. Kai kuris nors paleidimo modelis neturi
kainos, įvertis tai pasako (\code{pricesKnown} yra false), užuot pateikęs sumą, kuri tėra
apatinė riba.

\subsection{Eilutė}

\Cref{tab:models:llmcalls} išvardija visus eilutės laukus. Du iš jų verti paaiškinimo.

\code{kind} yra viena iš keturių reikšmių, kurias užtikrina duomenų bazės apribojimas:
\code{chat}, \code{job_step}, \code{embed}, \code{prompt_test}. \code{run_prompt} visada rašo
\code{job_step}, taip pat ir tiems keliems iškvietimams, kuriuos atlieka ne užduotyje (nuorodų
patikra iš teisės aktų bazės puslapio); tokia eilutė tiesiog neturi užduoties etapo.

\code{purpose} nurodo, kam iškvietimas buvo skirtas, ne daugiau kaip 80 simbolių.
Struktūrizuotam generavimui tai instrukcijos identifikatorius, nebent kviečiantysis nurodo kitą
(nuorodų vertintojas rašo \code{citations.check}; vertinimo rinkiniai rašo \texttt{eval:} ir
rinkinio pavadinimą, išskyrus nuorodų rinkinį, kurio iškvietimai eina per nuorodų vertintoją ir
neša jo paskirtį). Vektoriams tai viena iš keturių paskirčių (\cref{tab:models:embedcallers});
pokalbiui -- \code{chat} ėjimui ir \code{chat_title} pavadinimui.

\begin{xltabular}{\textwidth}{@{}L{38mm}Y@{}}
\caption{Lentelės \code{llm_calls} eilutės laukai (\file{assistant/models.py}, \code{LlmCall})}\label{tab:models:llmcalls}\\
\toprule
\hd{Laukas} & \hd{Saugo} \\
\midrule
\endfirsthead
\toprule
\hd{Laukas} & \hd{Saugo} \\
\midrule
\endhead
\bottomrule
\endlastfoot
\code{id} & Pačios eilutės identifikatorius. \\
\code{kind} & \code{chat}, \code{job_step}, \code{embed} arba \code{prompt_test}. \\
\code{purpose} & Kam iškvietimas buvo skirtas: instrukcijos identifikatorius arba viena iš aukščiau nurodytų fiksuotų paskirčių. \\
\code{user} & Naudotojas, kuriam priskiriama iškvietimo kaina; bendram darbui tuščias (\cref{tab:models:charged}). Dienos biudžetas sumuojamas pagal šį lauką. \\
\code{thread_id} & Pokalbio ėjimui -- pokalbis. \\
\code{job_step_id} & Foniniam darbui -- užduoties etapas. Etapui pasibaigus, užduočių vykdytojas etapo eilutes -- iškvietimus, žetonus, kainą, pagal modelį -- susumuoja į paties etapo įrašą (\file{jobs/runner.py}); iš ten ir imama ataskaitos kaina pagal etapus. \\
\code{prompt_version_id} & Instrukcijos versija, kurios instrukcijos buvo vykdytos; tuščias vektoriams ir neišsaugoto projekto bandomajam paleidimui. \\
\code{model} & Modelis, kuris buvo vykdomas, \texttt{provider:model} forma; iškvietimui per generavimo ar vektorizavimo maršrutą -- toks, kokį asistento paslauga nurodė savo atsakyme. \\
\code{input_tokens}, \code{cached_input_tokens}, \code{cache_write_tokens}, \code{output_tokens}, \code{reasoning_tokens} & Penki skaičiai. \\
\code{cost_usd} & Kaina, kurią \code{cost_of} apskaičiavo įrašant eilutę, šešių skaitmenų po kablelio tikslumu. Vėlesnis kainų pasikeitimas jos nekeičia. \\
\code{duration_ms} & Kiek truko iškvietimas. \\
\code{tool_calls} & Pokalbio ėjimui -- modelio iškviestų įrankių sąrašas. \\
\code{outcome} & \code{ok} arba, pokalbio ėjimui, kaip jis kitaip baigėsi (pvz., \code{timeout}, \code{budget}, \code{error}); pavadinimui -- \code{empty}, kai modelis jo neparašė. \\
\code{transcript} & Pokalbio ėjimui: naudotojo žinutė (daugiausia 20\,000 simbolių) ir asistento atsakymas (daugiausia 100\,000); pavadinimui -- pirmosios žinutės pradžia ir parašytas pavadinimas. Visiems kitiems iškvietimams tuščias. \\
\code{created_at} & Kada eilutė įrašyta (UTC). \\
\end{xltabular}

\paragraph{Stenograma.} Pokalbis yra vienintelė funkcija, kurios iškvietimai saugo savo turinį
sąnaudų eilutėje: ėjimui -- ko naudotojas klausė ir ką asistentas atsakė; pavadinimui -- ką
modelis perskaitė ir parašė. Duomenų modelio antraštinis komentaras tai vadina \enquote{auditu, ką
asistentas kam pasakė}. Ji pateikiama tik iškvietimo išsamiame rodinyje
(\route{GET}{/api/usage/calls/{id}}), kurį naudotojas gali skaityti savo iškvietimams, o
administratorius -- bet kuriems. Konvejerio iškvietimų įvestys ir išvestys į \code{llm_calls}
nekopijuojamos; jų rezultatai saugomi kiekvieno kelio lentelėse.

\paragraph{Telemetrija niekada nesugadina kviečiančiojo darbo.} \code{record_call} įterpimą
vykdo atskirame išsaugojimo taške ir sugauna kiekvieną išimtį: nepavykęs eilutės įrašymas
užrašomas žurnale, o modelio atsakymas vis tiek naudojamas.

\paragraph{Kas neužrašoma.} Du dalykai eilutės nepalieka. Pirmasis yra sąmoningas ir nurodytas
\file{docs/markdown/en/architecture.md}: paieškos užklausos vektorius, keli žetonai per paiešką,
sukuriamas be sąnaudų eilutės. Antrasis išplaukia iš kodo tvarkos: \code{run_prompt} ir
vektorių kviečiantieji iškvietimą užrašo tik gavę atsakymą, todėl nepavykęs užduočių vykdytojo
iškvietimas -- įskaitant tą, kurio atsakymą atmetė schemos patikra, -- eilutės neturi, nors
tiekėjas jo žetonus galėjo apmokestinti. Nepavykęs pokalbio ėjimas užrašomas su savo
\code{outcome}.

Konsolės sąnaudų puslapis šias eilutes skaito per \route{GET}{/api/usage} ir
\route{GET}{/api/usage/calls}: kaina per dieną, pagal rūšį, pagal modelį ir --
administratoriui -- pagal naudotoją (\cref{ch:console,ch:api}).


\section{Biudžetai}\label{sec:models:budgets}

Ką sistema gali išleisti, riboja du biudžetai. Vienas priklauso DI tinklų sąsajai ir apima visą
diegimą; kitas yra pačios sistemos ir apima vieną naudotoją vienai dienai.

\subsection{DI tinklų sąsajos rakto komandos biudžetas}

Raktas, kuriuo asistento paslauga kviečia DI tinklų sąsają, turi \term{komandos biudžetą},
kurį nustato ir skaičiuoja DI tinklų sąsaja. Niekas šioje sistemoje jo pakeisti negali. Kai jis
išnaudojamas, DI tinklų sąsaja į kiekvieną iškvietimą atsako HTTP~402, o sistema tai pasako
aiškiai, užuot pranešusi apie bendrą nesėkmę: asistento paslauga atsisakymą paverčia
402~\code{AI_BUDGET_EXHAUSTED} vidinėse duryse ir tuo pačiu kodu pokalbyje; taip nepavykusi
ataskaita užrašo pranešimą, prasidedantį tuo kodu, ir išsaugo viską, už ką jau buvo sumokėjusi
-- dokumentų inventorių ir išskirtų pareigų kitas paleidimas nebeperka iš naujo.

Dvi klasifikacijos (\code{classifyProviderError} faile \file{middleware/errors.js}) detalės
vertos dėmesio. 402 visada reiškia išnaudotą biudžetą; 400, 403, 429 ar klaida be būsenos
laikoma tokia tik tada, kai jos tekste yra ir pinigų žodis (budget, quota, billing, credit
balance), ir pasibaigimo žodis. Be to, dažnio limito formuluotė niekada nelaikoma išnaudotu
biudžetu, net jei mini kvotą: abu atvejai reikalauja priešingų reakcijų -- laukti arba sustoti.
\code{finalBudgetRefusals} tada biudžeto atsisakymą pažymi kaip nekartotiną, nes kitaip SDK
kartotų 429, kurio nepašalins joks kelių sekundžių laukimas.

\file{docs/markdown/en/operations.md} užrašo, kad komandos biudžetas 2026 m. rugsėjo 27 d., tą dieną,
kai buvo pasiektas, buvo 10~\$.

\subsection{Naudotojo dienos biudžetas}

Pačios sistemos biudžetas \enquote{tik neleidžia vienai paskyrai ištuštinti bendro rakto}, kaip
sakoma \file{docs/markdown/en/operations.md}. Kiekvienas naudotojas per UTC parą gali išleisti
nustatytą sumą:

\begin{itemize}
  \item numatytoji suma -- parametras \code{LLM_DAILY_BUDGET_USD}, nuo 2026~m. lapkričio 20~d.
    10,00~\$. Anksčiau ji buvo 2,00~\$, nustatyta rugsėjį, kai pokalbio ėjimas kainavo centus.
    Nuo lapkričio naudojamu pokalbio modeliu ėjimas, kuriuo perskaitomi įmonės dokumentai,
    kainuoja apie 2,90~\$, o vėlesnis juos perskaičiusio pokalbio ėjimas -- apie 1,00~\$
    (veikiančios sistemos lapkričio 17~d. ėjimų mediana buvo 1,02~\$), todėl pirmas dokumentų
    perdavimas išnaudodavo įprastos paskyros dieną, o kitas jos klausimas būdavo atmetamas;
  \item administratorius paskyrai gali nustatyti savą sumą (naudotojo \code{daily_budget_usd};
    konsolėje -- puslapyje \ui{Naudotojai}): tuščia reiškia numatytąją, 0 reiškia neribotą,
    kitaip -- suma iki 10\,000~\$;
  \item paties administratoriaus paskyra neribota, nebent jai nustatyta suma.
\end{itemize}

\code{budget_status(user_id)} faile \file{assistant/usage.py} apskaičiuoja būklę: limitą,
panaudotą sumą -- \code{cost_usd} sumą per naudotojo eilutes lentelėje \code{llm_calls} nuo
00:00~UTC -- ir likutį (nieko, kai limitas 0, t.~y. neribotas). Kadangi biudžetas sumuojamas
per sąnaudų eilutes, jis skaičiuoja lygiai tai, ką \cref{tab:models:charged} nurodo kaip
priskiriamą naudotojui, ir nieko daugiau.

Biudžetas tikrinamas \emph{prieš} pradedant darbą, niekada ne jo viduryje
(\cref{tab:models:budgetchecks}). Kad būtų pradėta ataskaita ar programos dokumentas, turi būti
likę 0,25~\$ (\code{MIN_BUDGET_TO_START_USD}); pradėtas paleidimas niekada nesustabdomas
pusiaukelėje, net jei viršija limitą. Apie paleidimo dydį naudotojui pasako įvertis, kurį jis
patvirtina prieš pradžią (\cref{ch:report}).

\begin{table}[H]
\caption{Kur tikrinamas dienos biudžetas ir kas atsitinka, kai jis išnaudotas}\label{tab:models:budgetchecks}
\begin{tabularx}{\textwidth}{@{}L{31mm}L{50mm}Y@{}}
\toprule
\hd{Funkcija} & \hd{Kur tikrinama} & \hd{Kai biudžetas išnaudotas} \\
\midrule
Pokalbio ėjimas & asistento paslauga, \file{routes/chat.js}, \code{checkBudget}, prieš kiekvieną ėjimą (klausia API per \code{/internal/assistant/budget}) & Kai nieko nelikę, ėjimas atmetamas su 429~\code{BUDGET_EXCEEDED}. Jei nepavyksta pati patikra, ėjimas vis tiek leidžiamas. \\
Ataskaitos pradžia & \file{analysis/api/views.py}, \route{POST}{/api/reports} & Atmetama su 429~\code{BUDGET_EXCEEDED}, kai likę mažiau nei 0,25~\$. \\
Ataskaitos apimties patvirtinimas & \file{analysis/api/views.py}, \route{POST}{/api/reports/{id}/scope} & Tas pats: brangiajai antrajai daliai reikia likusių 0,25~\$. \\
Programos dokumento pradžia & \file{programme/api/views.py}, \route{POST}{/api/companies/{id}/programme} & Tas pats. \\
Bandomasis paleidimas instrukcijų redaktoriuje & \file{prompts/api/views.py}, \code{test_run} & Kai nieko nelikę, atmetama su 429~\code{BUDGET_EXCEEDED}. \\
Nuorodų patikra iš teisės aktų bazės puslapio ar pokalbio & \file{corpus/api/views.py}, \route{POST}{/api/corpus/check-citations} & Neatmetama: patikra vykdoma, vertinant pagal žodžių sutapimą, o ne modeliu. \\
Įkelto dokumento indeksavimas & \file{companies/indexing.py}, užduotis \code{documents.index} & Niekas nevektorizuojama. Užduotis pati save vėl įrašo į eilę kitos dienos 00:05~UTC ir praneša, kad laukia biudžeto. \\
\bottomrule
\end{tabularx}
\end{table}

Paskutinė eilutė turi savo priežastį, nurodytą kode: be laukimo dokumentą cikle įkeliant ir
ištrinant, jis būtų vektorizuojamas be galo, o DI tinklų sąsajai būtų mokama be paliovos.

\subsection{Kam kas priskiriama}

Ar iškvietimas skaičiuojamas naudotojui, lemia vienas argumentas: \code{user_id}, kurį
kviečiantysis perduoda \code{run_prompt} ar \code{record_call}. Pasirinkimų taisyklė tokia:
darbas, atliktas \emph{vieno naudotojo įmonei}, priskiriamas tam naudotojui, o darbas,
atliktas \emph{su bendra teisės aktų baze} -- kurią pakartotinai naudoja kiekviena vėlesnė
kiekvienos įmonės ataskaita, -- niekam (\cref{tab:models:charged}). Bendras darbas vis tiek
naudoja DI tinklų sąsajos rakto komandos biudžetą; jis tik neskaičiuojamas į žmogaus dienos
biudžetą.

\begin{table}[H]
\caption{Kam priskiriama modelio iškvietimo kaina}\label{tab:models:charged}
\begin{tabularx}{\textwidth}{@{}L{62mm}Y@{}}
\toprule
\hd{Darbas} & \hd{Priskiriama} \\
\midrule
Pokalbio ėjimas ir naujam pokalbiui parašytas pavadinimas & Besikalbančiam naudotojui. \\
Ataskaita: dokumentų inventoriai, teisės aktų nuorodų patikra, taikytinumas, atitikties spragų analizė, veiksmų planas ir santrauka; ataskaitos aptiktų neindeksuotų dokumentų indeksavimas & Naudotojui, kuris pradėjo ataskaitą; po apimties peržiūros antrosios dalies iškvietimai priskiriami tam, kas patvirtino apimtį (tos dalies užduotis yra jo). \\
Įmonės vidaus reikalavimai, perskaityti iš jos politikų ir įvertinti ataskaitoje & Naudotojui, kuris pradėjo ataskaitą. \\
Programos dokumentas & Naudotojui, kuris jį pradėjo, išskyrus nuorodų vertintojo iškvietimą jo nuorodų etape: jis užrašomas prie užduoties etapo be naudotojo. \\
Įkelto dokumento indeksavimas (vektoriai) & Įkėlusiajam. \\
Nuorodų patikra iš teisės aktų bazės puslapio ar pokalbio & Naudotojui, kuris paprašė. \\
Bandomasis paleidimas instrukcijų redaktoriuje & Administratoriui, kuris jį paleido. \\
Pareigų išskyrimas, kurį ataskaita sukelia, nes aktas niekada nebuvo perskaitytas & Niekam: bendras teisės aktų bazės darbas. \\
Pareigų išskyrimas arba akto pareigų vertintojas, į eilę įrašyti administratoriaus & Tam administratoriui. \\
Akto apimties ir jo sankcijų nuostatų skaitymas; standarto struktūros skaitymas & Niekam. \\
Akto įkėlimas (jo fragmentų vektoriai); pareigų vektoriai & Niekam. \\
Paieškos užklausos vektorius & Visai neužrašoma. \\
\bottomrule
\end{tabularx}
\end{table}

\begin{humanbox}{Ką sprendžia administratorius}
Jokio biudžeto sistema pati nenustato ir nepanaikina. Administratorius sprendžia kiekvienos
paskyros dienos sumą ir gali ją panaikinti; DI tinklų sąsajos operatoriai sprendžia komandos
biudžetą. Administratorius taip pat kiekvienai instrukcijai sprendžia, kuris modelis vykdo
(\cref{sec:models:defaults}), -- tai didžiausias atskiras svertas, lemiantis, kiek kainuoja
ataskaita. Naudotojas nusprendžia pradėti paleidimą, pamatęs jo kainos įvertį.
\end{humanbox}


\section{Instrukcijų podėlis}\label{sec:models:cache}

Už DI tinklų sąsajos esantis tiekėjas laiko \term{instrukcijų podėlį} (angl. \en{prompt
cache}): kai užklausa prasideda tekstu, kurį jis jau apdorojo, tą pradžią jis nuskaito iš
podėlio ir apmokestina nedidele įvesties kainos dalimi. Sistema to nesirinko, tačiau jos
sprendimai turėjo prie to prisitaikyti, nes didžiausi ataskaitos iškvietimai didžiąją dalį
savo įvesties kartoja.

\subsection{Kas laikoma podėlyje, o kas ne}

Ką tiksliai tiekėjas laiko podėlyje, nustatyta matuojant, ir rezultatas siauresnis nei
\enquote{viskas, kas kartojasi}.

\begin{lessonbox}{Išmatuota 2026 m. lapkričio 10--11 d., gpt-6.1-sol}
Tas pats 8\,800 žetonų blokas buvo išsiųstas du kartus iš eilės keliose pozicijose. Antrą
kartą jis buvo \textbf{nuskaitytas iš podėlio}, kai stovėjo užklausos instrukcijose. Jis buvo
\textbf{įrašytas iš naujo} -- apmokestintas visa kaina, -- kai stovėjo naudotojo žinutės,
kurios pabaiga skyrėsi, pradžioje; kai buvo atskira žinutė prieš kitokią žinutę; kai su juo
buvo išsiųstas podėlio raktas; ir kai su juo buvo išsiųsta saugojimo trukmės užuomina.
Užklausa, kuri nuo pat pradžios žodis į žodį kartojo ar pratęsė ankstesnę užklausą, tai, ką
kartojo, nuskaitė iš podėlio.

Taigi tiekėjas podėlyje laiko du dalykus: užklausos \emph{instrukcijas} ir \emph{visą
ankstesnę užklausą}, kurią vėlesnė kartoja ar pratęsia. Iš žinutės, kuri tik prasideda kaip
ankstesnė, podėlyje nelieka nieko.

Tos dienos \model{gpt-6.1-sol} kainos: iš podėlio nuskaityta įvestis -- 0,10~\$ už milijoną
žetonų, dvidešimtoji įvesties kainos dalis; kiekvienas kitas įvesties žetonas apmokestinamas
kaip įrašytas į podėlį -- 2,50~\$, ketvirčiu daugiau nei 2~\$ įvesties kaina.
\end{lessonbox}

\subsection{Trys pasekmės sprendimams}

\paragraph{1. Tai, kas bendra etapo iškvietimams, dedama po instrukcijų.} Ataskaitos atitikties
spragų analizė su kiekvienu iškvietimu siunčia įmonės dokumentus. Jei jie būtų likę naudotojo
žinutėje, prieš pareigas, kurios kiekviename iškvietime skiriasi, kiekvienas iškvietimas už
juos būtų mokėjęs visą kainą. Todėl jie siunčiami atskirai, kaip \code{shared}, o generavimo
maršrutas juos prijungia prie instrukcijų -- vienintelės vietos, kurią tiekėjas laiko
podėlyje (\cref{fig:models:anatomy}).

\paragraph{2. Pirmasis tokio etapo iškvietimas vykdomas vienas.} Podėlis bloką laiko tik tada,
kai apdorotas iškvietimas, kuris jį nešė. \code{shared_first} vykdo vieną iškvietimą, jo
palaukia ir tik tada pradeda kitus po keturis vienu metu. Juo naudojasi ataskaitos
taikytinumo etapas ir atitikties spragų analizė, taip pat įmonės vidaus reikalavimų vertinimas.

\paragraph{3. Sąnaudos užrašomos ir įkainojamos dalis po dalies.} Tai ir yra
penkių skaičių (\cref{sec:models:usage}) priežastis. Iki 2026 m. lapkričio 11 d. kiekvienas
įvesties žetonas buvo įkainojamas paprasta įvesties kaina; \file{docs/markdown/en/architecture.md}
užrašo, kad pokalbio ėjimas, kuris kiekviename savo žingsnyje vėl siunčia instrukcijas, tada
buvo užrašomas iki dviejų kartų didesne nei tikroji kaina. Biudžetai ir ataskaitos kaina pagal
etapus remiasi pataisytais skaičiais.

\begin{figure}[H]
\centering
\begin{tikzpicture}[x=1mm, y=1mm]
  \tikzset{seg/.style={draw=ink, line width=.6pt, minimum height=9mm, align=center,
                       inner sep=2pt, font=\sffamily\footnotesize, text=body}}
  % the parts of the request
  \node[seg, fill=paper, minimum width=24mm, anchor=west] (i0) at (26,0) {instrukcijos};
  \node[seg, fill=paper, minimum width=66mm, anchor=west] (s0) at (50,0) {bendrasis blokas: profilis ir dokumentai};
  \node[seg, fill=paper, minimum width=30mm, anchor=west] (d0) at (116,0) {iškvietimo duomenys};
  \node[note, anchor=east] at (24,0) {viena užklausa};
  \draw[decorate, decoration={brace, amplitude=4pt}, draw=muted] (26,6.5) -- node[note, above=4pt, align=center] {siunčiama kaip instrukcijos -- dalis, kurią tiekėjas laiko podėlyje} (116,6.5);
  \draw[decorate, decoration={brace, amplitude=4pt}, draw=muted] (116,6.5) -- node[note, above=4pt, align=center] {naudotojo žinutė} (146,6.5);
  % the first call
  \node[seg, fill=warnbg, minimum width=120mm, anchor=west] (a) at (26,-17) {įrašyta į podėlį: apie 59\,000 žetonų po 2,50~\$ už milijoną};
  \node[note, anchor=east, align=right] at (24,-17) {pirmasis\\iškvietimas,\\vienas};
  % later calls
  \node[seg, fill=goodbg, minimum width=90mm, anchor=west] (b1) at (26,-31) {nuskaityta iš podėlio: apie 55\,000 žetonų po 0,10~\$ už milijoną};
  \node[seg, fill=warnbg, minimum width=30mm, anchor=west] (b2) at (116,-31) {įrašyta};
  \node[note, anchor=east, align=right] at (24,-31) {kiekvienas\\vėlesnis\\iškvietimas};
\end{tikzpicture}
\caption{Užklausos, nešančios bendrąjį bloką, sandara, su pavyzdinės įmonės atitikties spragų
analizės skaičiais (26 politikos, 261\,000 simbolių; išmatuota 2026 m. lapkričio 10--11 d.).
Instrukcijos ir bendrasis blokas kiekviename etapo iškvietime yra vienodi ir nuskaitomi iš
tiekėjo podėlio, kai pirmasis iškvietimas juos įrašė; kiekvieno iškvietimo duomenys kaskart
yra nauja įvestis.}\label{fig:models:anatomy}
\end{figure}

\subsection{Išmatuotas poveikis}

Pavyzdinėje įmonėje kiekvienas atitikties spragų analizės iškvietimas neša apie 59\,000
įvesties žetonų, beveik visi jie -- 26 įmonės politikos. \Cref{fig:models:cachecost} rodo, kiek
vienas iškvietimas kainavo trimis atvejais, kaip juos užrašo \file{docs/markdown/en/architecture.md}
(2026 m. lapkričio 10--11 d., \model{gpt-6.1-sol}).

\begin{figure}[H]
\centering
\begin{tikzpicture}
\begin{axis}[
  xbar, width=92mm, height=50mm,
  xmin=0, xmax=0.22, bar width=6.5mm,
  enlarge y limits=0.28,
  axis x line*=bottom, axis y line*=left,
  axis line style={draw=muted}, tick style={draw=muted},
  ytick=data, symbolic y coords={before,read,write},
  yticklabels={{dokumentai naudotojo žinutėje\\(iki pakeitimo)},
               {vėlesnis iškvietimas: blokas\\nuskaitytas iš podėlio},
               {pirmasis iškvietimas: blokas\\įrašytas į podėlį}},
  yticklabel style={align=right, text width=50mm, font=\sffamily\footnotesize, text=body},
  xtick={0,0.05,0.10,0.15,0.20},
  xticklabel style={font=\sffamily\footnotesize, text=muted, /pgf/number format/fixed,
                    /pgf/number format/precision=2, /pgf/number format/fixed zerofill,
                    /pgf/number format/use comma},
  scaled x ticks=false,
  xlabel={JAV doleriai už vieną iškvietimą}, xlabel style={font=\sffamily\footnotesize, text=muted},
  nodes near coords={\pgfmathprintnumber[fixed, precision=3, use comma]{\pgfplotspointmeta}~\$},
  nodes near coords style={font=\sffamily\footnotesize, text=body, anchor=west},
  point meta=x]
  \addplot[fill=mid, draw=ink] coordinates {(0.17,before) (0.045,read) (0.18,write)};
\end{axis}
\end{tikzpicture}
\caption{Vieno pavyzdinės įmonės atitikties spragų analizės iškvietimo kaina (apie 59\,000
įvesties žetonų, \model{gpt-6.1-sol}, išmatuota 2026 m. lapkričio 10--11 d.). Pirmasis etapo
iškvietimas sumoka 0,18~\$ už bendrojo bloko įrašymą; kiekvienas vėlesnis iškvietimas 55\,000
jo žetonų nuskaito iš podėlio ir kainuoja apie 0,045~\$. Iki pakeitimo, kai dokumentai buvo
naudotojo žinutėje, kiekvienas iškvietimas kainavo 0,17~\$.}\label{fig:models:cachecost}
\end{figure}

\file{docs/markdown/en/architecture.md}, remdamasis 2026 m. lapkričio 11 d. paleidimu duomenų bazės
kopijoje, kurioje kiekvieno akto pareigos buvo apkarpytos iki pirmųjų penkiolikos, ekstrapoliuoja,
kad visa ataskaita tada būtų pareikalavusi apie 82 tokių iškvietimų ir kainavusi maždaug 5~\$, palyginti su
16,04~\$ už dieną anksčiau parengtą ataskaitą; prie šio skirtumo prisideda ir kiti konvejerio
pakeitimai tarp abiejų paleidimų (\cref{ch:report}). Pirmoji visa ataskaita veikiančioje
sistemoje, 2026 m. lapkričio 13 d., su iki 68 įjungtų aktų išaugusia teisės aktų baze, atliko
154 atitikties spragų analizės iškvietimus: jie išsiuntė 9,12 milijono įvesties žetonų, iš
kurių 8,44 milijono (92,5~\%) buvo nuskaityti iš tiekėjo podėlio, o 0,68 milijono į jį
įrašyti. Tas pats mechanizmas be jokio papildomo sprendimo tarnauja ir pokalbiui: kiekvienas
ėjimo žingsnis pratęsia ankstesnio žingsnio užklausą, todėl žingsnis savo instrukcijas ir
ankstesnius žingsnius nuskaito iš podėlio, o visą kainą moka tik už tai, ką pridėjo paskutinis
žingsnis (\cref{ch:chat}).

Podėlis priklauso tiekėjui, ir sistemos teisingumas nuo jo nepriklauso. Iškvietimas, radęs
tuščią podėlį, grąžina tą patį atsakymą ir kainuoja daugiau; sąnaudų eilutė rodo, kuris iš šių
dviejų atvejų įvyko.


\section{Dažnio limitai}\label{sec:models:ratelimits}

Be pinigų, tiekėjas riboja pralaidumą: jis skaičiuoja, kiek žetonų raktas išsiunčia per
minutę. Daugybė didelių ataskaitos iškvietimų tą limitą gali pasiekti per kelias minutes.

\begin{lessonbox}{Nesėkmė, 2026 m. lapkričio 10 d.}
Keturi atitikties spragų analizės iškvietimai po 59\,000 žetonų, vykdomi vienu metu, per dvi
minutes pasiekė tiekėjo žetonų per minutę limitą. Trumpi kliento pakartojimai -- po 2, 6 ir
18~sekundžių -- vėl atsimušė į limitą, ir ataskaita nepavyko. Nuo tada užduoties etapui
atliekamas iškvietimas dažnio limito išlaukia -- tiek, kiek prašo tiekėjas, ir iš viso iki
penkiolikos minučių (\code{RATE_LIMIT_PATIENCE_S}), -- užuot paversdamas savo etapą
nesėkmingu.
\end{lessonbox}

\paragraph{Kaip ateina dažnio limitas.} Jis ateina viena iš dviejų formų, ir
\code{providerPause} faile \file{middleware/errors.js} skaito abi. Įprasta forma -- HTTP 429,
kuris nėra biudžeto atsisakymas, su antrašte \code{Retry-After} arba be jos. Kita forma buvo ta,
su kuria susidurta lapkričio 10 d.: tiekėjas jau buvo pradėjęs savo srautą ir apie limitą
pranešė kaip apie klaidos įvykį jo viduje, visai be HTTP būsenos; SDK tai perduoda kaip srauto
analizės nesėkmę. Tokia klaida atpažįstama iš jos žodžių (\enquote{rate limit},
\enquote{call rate}, \enquote{retry after} ir skaičius). Laukimas imamas iš antraštės, jei ji
yra, kitaip -- iš pranešime nurodytų sekundžių, kitaip jis nežinomas.

\paragraph{Kas su juo daroma.} SDK asistento paslaugoje pirmiausia du kartus pakartoja
iškvietimą. Limitas, kuris trunka ilgiau, arba limitas, atėjęs srauto viduje, API perduodamas
kaip 429 su lauku \code{retryAfterS}, jei tiekėjas nurodė laukimą. Nuo čia taikomos
taisyklės (\cref{tab:models:retries}): užduoties etapas laukia ir klausia dar kartą, nelaikydamas
to nesėkme; iškvietimas, kurio laukia žmogus, trumpai kartojamas, o tada nepavyksta su
pranešimu. Pokalbyje dažnio limitas rodomas kaip pranešimas pabandyti dar kartą
(\cref{ch:chat}).

Ataskaitos etapas vienu metu vykdo keturis iškvietimus (\code{WORKERS} faile
\file{analysis/pipeline.py}); limitą sugeria kiekvieno iškvietimo kantrybė.

\paragraph{Pačios sistemos pokalbio ėjimų limitas.} Vienas dažnio limitas yra pačios sistemos
ir susijęs su žmonėmis, ne su žetonais: naudotojas per minutę gali pradėti daugiausia 12
pokalbio ėjimų (\code{RATE_LIMIT_TURNS}, fiksuotas \code{RATE_LIMIT_WINDOW_MS} = 60\,000~ms
langas, kurį atmintyje laiko \file{services/ratelimit.js}). Limitą viršijantis ėjimas
atmetamas su 429~\code{RATE_LIMITED} ir lange likusių sekundžių skaičiumi (\cref{ch:chat}).


\section{Kiekvieno etapo numatytasis modelis}\label{sec:models:defaults}

Kuriuo modeliu vykdoma instrukcija, yra jos aktyvios versijos laukas, ir administratorius jį
gali pakeisti instrukcijų redaktoriuje (\cref{ch:prompts}). Su sistema pristatomos versijos
nurodo tris kalbos modelius ir vieną vektorizavimo modelį. \Cref{tab:models:defaults} pateikia
kiekvienos instrukcijos numatytąjį modelį ir užrašytą jo priežastį. Kiekviena priežastis -- tai
datuotas stebėjimas, atliktas su pavyzdinės įmonės medžiaga; nė viena nėra bendras modelių
reitingas.

\begin{xltabular}{\textwidth}{@{}L{36mm}L{40mm}Y@{}}
\caption{Kiekvienos instrukcijos numatytasis modelis su užrašyta priežastimi. Kainos -- JAV doleriai už milijoną įvesties ir išvesties žetonų.}\label{tab:models:defaults}\\
\toprule
\hd{Instrukcijos} & \hd{Numatytasis modelis} & \hd{Kodėl} \\
\midrule
\endfirsthead
\toprule
\hd{Instrukcijos} & \hd{Numatytasis modelis} & \hd{Kodėl} \\
\midrule
\endhead
\bottomrule
\endlastfoot
\prompt{chat.assistant} & \model{gpt-6-astra} (10 / 50) & Pristatomas pokalbio numatytasis modelis; prieš jį buvo \model{gpt-5.5}. Saugyklos dokumentuose pokalbiui jokio palyginimo neužrašyta. Naudotojas pokalbiui gali pasirinkti kitą modelį (\cref{ch:chat}). \\
\prompt{chat.title} & \model{gpt-5.6-luna} (0,20 / 1,20) & Keli žodžiai, parašyti pagal pirmosios pokalbio žinutės pradžią: mažasis modelis. \\
Ataskaita: \prompt{report.inventory}, \prompt{report.applicability}, \prompt{report.gap_analysis}, \prompt{report.action_plan}, \prompt{report.plan_merge}, \prompt{report.summary}; įmonės vidaus reikalavimai: \prompt{company.internal_requirements}, \prompt{report.internal_gap} & \model{gpt-6.1-sol} (2 / 10) & Išmatuota 2026 m. lapkričio 10 d. su tikromis pavyzdinės įmonės užklausomis -- trimis atitikties spragų analizės paketais po penkiolika pareigų ir dviem taikytinumo paketais po keturis aktus: \model{gpt-6.1-sol} ir \model{gpt-6-astra} 43 iš 45 pareigų nustatė tą pačią būseną ir visiems aštuoniems aktams -- tą patį verdiktą, už penktadalį kainos; \model{gpt-5.6-terra} skyrėsi penkiose pareigose, kiekvieną kartą praleisdamas ištrauką arba reikalaudamas iš politikos to, ko pareiga nereikalauja. Veiksmų planas ir santrauka, kiekvieno modelio parašyti iš tų pačių dvidešimties atvirų išvadų, apėmė visas dvidešimt. \\
\prompt{corpus.extract_obligations}, \prompt{corpus.extract_requirements} & \model{gpt-6.1-sol} & Vienkartinis teisės aktų bazės darbas, kurį pakartotinai naudoja kiekviena ataskaita, todėl -- kruopščiausias išmatuotas skaitytojas. 2026 m. lapkričio 10 d., dviem paketais (MiCA 68--75 str. ir Pinigų plovimo ir teroristų finansavimo prevencijos įstatymo, PPTFPĮ, 10--13 str.): \model{gpt-5.6-terra} išvardijo 44 ir 25 pareigas, \model{gpt-6.1-sol} -- 56 ir 32, \model{gpt-6-astra} -- 55 ir 31, maždaug penkis kartus brangiau. \\
\prompt{corpus.standard_outline}, \prompt{corpus.scope}, \prompt{corpus.sanctions} & \model{gpt-6.1-sol} & Taip pat vienkartinis teisės aktų bazės darbas; \file{docs/markdown/en/architecture.md} standarto struktūrą ir akto apimtį priskiria tai pačiai grupei kaip išskyrimą. Šioms trims savo palyginimo neužrašyta. \\
\prompt{corpus.check_citations} & \model{gpt-6.1-sol} & Trumpi palyginimai, keli iškvietimai ataskaitai. Patikra pirmiausia buvo išmatuota su \model{gpt-5.6-luna}: jis teisingai įvertino visas 12 klaidingų pavyzdinės įmonės nuorodų į MiCA, o vektorių panašumas buvo ne geresnis už atsitiktinumą. \\
\prompt{corpus.verify_obligations} (vertintojas) & \model{gpt-6.1-sol} & Vienkartinis teisės aktų bazės darbas: viso 5\,960 pareigų registro įvertinimas 2026 m. lapkričio 13 d. pareikalavo 369 iškvietimų ir 9,80~\$. \\
Šešios programos dokumentų instrukcijos, kurių identifikatoriai prasideda \prompt{programme.} & \model{gpt-6.1-sol} & Du iškvietimai vienam dokumentui per visą įmonės dokumentų rinkinį; antrasis rinkinį nuskaito iš podėlio. Dokumentas kainuoja apie 0,4--0,5~\$ (\file{docs/markdown/en/architecture.md}). \\
Vektoriai (be instrukcijos) & \model{text-embedding-3-large}, 1536 matmenys (0,13) & 2026 m. rugsėjo 28 d. palyginimas (\cref{tab:models:embedcompare}). \\
\end{xltabular}

\paragraph{Vienas modelis visai ataskaitai.} \file{docs/markdown/en/architecture.md} iš lapkričio 10 d.
matavimo daro išvadą: \model{gpt-6-astra} kainuoja penkis kartus daugiau nei
\model{gpt-6.1-sol} už sutapusius atsakymus -- apie 57~\$ palyginti su 11~\$ už pavyzdinės
įmonės 2\,200 pareigų atitikties spragų analizę, -- todėl kiekvienas ataskaitos etapas
vykdomas \model{gpt-6.1-sol}. Tas pats dokumentas nurodo ir šio pasirinkimo kainą laiku: tą
dieną GPT-6 modeliai rašė apie 70 žetonų per sekundę, perpus lėčiau nei
\model{gpt-5.6-terra}, todėl tankus paketas trunka nuo keturių iki šešių minučių, o
viršutinės ribos (\cref{tab:models:timeouts}) yra tokios, kokios yra.

\paragraph{Vektorizavimo modelis.} 2026 m. rugsėjo 28 d. su pačios sistemos duomenimis buvo
palyginti du DI tinklų sąsajos modeliai (\cref{tab:models:embedcompare}). Didesnysis aiškiai
geresnis būtent ten, kur sistemai to reikia: ieškant lietuviškos nuostatos pagal anglišką tekstą.
Pokalbis dažnai ieško angliškai, įmonių dokumentai dažnai būna anglų kalba, o teisės aktų
bazė yra lietuviška. \file{docs/markdown/en/architecture.md} didesniojo modelio kainą nurodo apie
0,13~\$ už milijoną žetonų, maždaug 0,30~\$ už viso turinio pakartotinį vektorizavimą.

\begin{table}[H]
\caption{Du palyginti vektorizavimo modeliai, 2026 m. rugsėjo 28 d. (pagal \file{docs/markdown/en/architecture.md}). Pirmosios trys eilutės skaičiuoja, kiek kartų iš 40 pareigos tekstas tarp lietuviškų akto fragmentų rado savo paties straipsnį.}\label{tab:models:embedcompare}
\begin{tabularx}{\textwidth}{@{}Yrr@{}}
\toprule
\hd{Bandymas} & \hd{-large} & \hd{-small} \\
\midrule
Angliškas pareigos tekstas, MiCA & 38 & 23 \\
Angliškas pareigos tekstas, DORA & 33 & 29 \\
Angliškas pareigos tekstas, PPTFPĮ & 37 & 34 \\
Lietuviškos užklausos (MiCA / DORA / PPTFPĮ) & 39 / 35 / 37 & 36 / 37 / 36 \\
Įmonės įrodymai: tinkamas dokumentas tarp šešių artimiausių pareigos ištraukų, iš 41 & 38 & 35 \\
Tas pats, vidutinis atvirkštinis rangas & 0,83 & 0,79 \\
\bottomrule
\end{tabularx}
\end{table}

Ką modeliams sakoma kiekvienoje iš šių instrukcijų ir kaip jų atsakymai tikrinami ne tik pagal formą,
aptariama tolesniuose skyriuose; \cref{app:prompts} išvardija kiekvieną instrukciją, o
\cref{ch:control} surenka visai ataskaitai išmatuotas kainas.
